\documentclass[../main.tex]{subfiles}
\begin{document}

\chapter{メモリ管理}
\lessoninfo{
  video={P3L2},
  textbook={OSTEP \cite{ostep} ch.~13，15--23；Silberschatz ら \cite{silberschatz2018} ch.~9--10；Bovet と Cesati \cite{bovet2005} ch.~2，8，17},
  keywords={仮想メモリ，ページング，セグメンテーション，ページテーブルエントリ，ページフォールト，多段ページテーブル，TLB，ラージページ，バディアロケータ，スラブアロケータ，デマンドページング，ページ置換，コピーオンライト，チェックポインティング},
}

第7章では，オペレーティングシステムが CPU を複数のプロセスでどのように
分け合うかを述べた．この章では，あらゆるプロセスが必要とする第2の資源，
すなわち物理メモリを扱う．オペレーティングシステムは，各プロセスに専用の
仮想アドレス空間（第2章）を与え，使用中の部分だけを物理メモリに保持し，
他の部分は必要になったときに二次記憶から持ち込む．この対応付けを構成する
2つの方式であるページとセグメント，アドレス変換を高速かつコンパクトにする
ハードウェアとデータ構造，物理メモリを管理するアロケータ，どのページが
メモリを去るかを決める方針，そしてメモリ管理のハードウェアによって可能に
なる2つのサービス，すなわちコピーオンライトとチェックポインティングを順に
調べる．

\section{メモリ管理の目標}
\label{sec:l08-goals}

オペレーティングシステムは，マシンの物理メモリ（DRAM）を，その上で動作する
プロセスに代わって管理する．\source{P3L2，レッスンの概観，メモリ管理の
目標；OSTEP \cite{ostep} ch.~13，15；Silberschatz ら
\cite{silberschatz2018} ch.~9．}プロセスは物理アドレスを使わない．各
プロセスは仮想アドレス空間を持ち，その範囲は物理メモリよりはるかに大きく
なりうる．32 ビットの仮想アドレスでは，どのプロセスも \SI{4}{\gibi\byte}
をアドレス指定でき，すべてのプロセスのアドレス空間を合わせれば，搭載されて
いるメモリを容易に超える．

\begin{definition}[仮想メモリ]
  オペレーティングシステムが，プロセスの仮想アドレス空間を物理メモリから
  切り離すために用いる抽象化を\term[かそうめもり]{仮想メモリ}[virtual memory]
  と呼ぶ．オペレーティングシステムは，各アドレス空間のうち使用中の部分を
  物理メモリに対応付け，他の部分を二次記憶に置き，プロセスが使うすべての
  アドレスを変換し，検証する．
\end{definition}

プロセスは自分の状態のすべてを同時に必要とするわけではない．どの時点でも
アドレス空間の一部しか使わない．したがってオペレーティングシステムは，
プロセスが次に必要とする状態を速やかに利用可能にできる限り，その部分だけを
物理メモリに保持すればよい．ここから2つの仕事が生じる．
\begin{description}
  \item[割り当て] オペレーティングシステムは，どのアドレス空間のどの部分を
    物理メモリのどこに置くかを決め，仮想アドレスから物理アドレスへの
    対応付けを確立する．物理メモリはプロセスが使うメモリより小さいことが
    あるので，一部の状態はディスクへ移し，必要になったときに戻さなければ
    ならない．したがって割り当てには\emph{置換}，すなわち場所を空けるために
    メモリを去る状態の選択が含まれる．
  \item[調停] すべてのメモリアクセスは，仮想アドレスから物理アドレスへ
    変換され，かつ検証されなければならない．そのアドレスはアドレス空間の
    割り当て済みの領域に属していなければならず，そのアクセスは許可されて
    いなければならない．仮想アドレスを，それが対応付けられた物理アドレスへ
    変換することを\term[あどれすへんかん]{アドレス変換}[address translation]
    と呼ぶ．これはメモリアクセスのたびに起こるので，高速でなければならない．
\end{description}

オペレーティングシステムがアドレス空間を単位に分割する方法は2つある
（\cref{fig:l08-goals}）．
\begin{description}
  \item[ページに基づくメモリ管理] \term[ぺーじんぐ]{ページング}[paging]では，
    仮想アドレス空間は固定長のブロックに分割され，その各々を
    \term[ぺーじ]{ページ}[page]と呼び，物理メモリも同じ大きさのブロックに
    分割され，その各々を\term[ぺーじふれーむ]{ページフレーム}[page frame]と
    呼ぶ．割り当てはページをページフレームに対応付け，調停はページテーブルを
    用いる．
  \item[セグメントに基づくメモリ管理]
    \term[せぐめんてーしょん]{セグメンテーション}[segmentation]では，
    アドレス空間は可変長のブロックに分割され，その各々を
    \term[せぐめんと]{セグメント}[segment]と呼ぶ．セグメントは，コード，
    ヒープ，スタックのような，プロセスの論理的な構成要素に対応する．
    割り当てはセグメントを同じ大きさの物理メモリの領域に対応付け，調停は
    セグメントレジスタを用いる．
\end{description}
現在のオペレーティングシステムではページングが主流である．セグメンテーション
は一部のアーキテクチャ，とりわけ x86 でサポートされており，そこではページング
と組み合わされる（\cref{sec:l08-segmentation}）．

\begin{figure}[htbp]
  \centering
  % Page-based versus segment-based memory management: fixed-size pages
% mapped to page frames (one page on disk, one unmapped) and variable-size
% segments mapped to regions of physical memory.
% Usage: % Page-based versus segment-based memory management: fixed-size pages
% mapped to page frames (one page on disk, one unmapped) and variable-size
% segments mapped to regions of physical memory.
% Usage: % Page-based versus segment-based memory management: fixed-size pages
% mapped to page frames (one page on disk, one unmapped) and variable-size
% segments mapped to regions of physical memory.
% Usage: \input{figures/l08-goals}   (width ~118 mm)
\begin{tikzpicture}[x=1mm, y=1mm, font=\scriptsize\sffamily,
  >={Stealth[length=1.4mm]},
  pg/.style={draw, fill=ln-accent!22},
  sg/.style={draw, fill=ln-source!16},
  oth/.style={draw, fill=ln-box},
  lab/.style={font=\scriptsize\sffamily},
  hd/.style={font=\scriptsize\sffamily\bfseries, anchor=base},
  capt/.style={font=\footnotesize\sffamily}]
  % ================= (a) paging
  \node[hd] at (9,38) {\iflangja{仮想}{virtual}};
  \node[hd] at (45,38) {\iflangja{物理}{physical}};
  % virtual pages 0..5 (height 6)
  \foreach \j/\s in {0/pg, 1/pg, 2/pg, 4/pg, 5/pg} {
    \draw[\s] (2,6*\j) rectangle ++(14,6);
    \node[lab] at (9,3+6*\j) {\iflangja{ページ}{page} \j};
  }
  \draw[fill=white] (2,18) rectangle ++(14,6);
  \node[lab, text=ln-muted] at (9,21) {\iflangja{未使用}{unused}};
  % page frames 0..7 (height 4.5)
  \foreach \i in {0,...,7} { \draw[oth] (38,4.5*\i) rectangle ++(14,4.5); }
  \foreach \i in {0,2,5,7} { \draw[pg] (38,4.5*\i) rectangle ++(14,4.5); }
  \foreach \i in {1,3,4,6} { \node[lab, text=ln-muted] at (45,2.25+4.5*\i) {\iflangja{他}{other}}; }
  % disk
  \draw[fill=ln-box, rounded corners=2pt] (38,-11) rectangle (52,-4)
    node[midway, lab] {\iflangja{ディスク}{disk}};
  % mappings: page -> frame
  \foreach \j/\i in {0/5, 1/2, 4/7, 5/0} {
    \draw[->] (16,3+6*\j) -- (38,2.25+4.5*\i);
  }
  \draw[->] (16,15) -- (38,-7.5);
  \node[capt] at (27,-15) {\iflangja{(a) ページング}{(a) paging}};
  % ================= (b) segmentation
  \node[hd] at (72,38) {\iflangja{仮想}{virtual}};
  \node[hd] at (108,38) {\iflangja{物理}{physical}};
  \draw[sg] (65,0)  rectangle (79,9)  node[midway, lab] {\iflangja{コード}{code}};
  \draw[sg] (65,9)  rectangle (79,14) node[midway, lab] {\iflangja{データ}{data}};
  \draw[sg] (65,14) rectangle (79,22) node[midway, lab] {\iflangja{ヒープ}{heap}};
  \draw[fill=white] (65,22) rectangle (79,30);
  \draw[sg] (65,30) rectangle (79,36) node[midway, lab] {\iflangja{スタック}{stack}};
  % physical memory
  \draw[oth] (101,-11) rectangle (115,36);
  \draw[sg] (101,-9) rectangle (115,-3) node[midway, lab] {\iflangja{スタック}{stack}};
  \draw[sg] (101,2)  rectangle (115,11) node[midway, lab] {\iflangja{コード}{code}};
  \draw[sg] (101,16) rectangle (115,24) node[midway, lab] {\iflangja{ヒープ}{heap}};
  \draw[sg] (101,28) rectangle (115,33) node[midway, lab] {\iflangja{データ}{data}};
  \draw[->] (79,4.5) -- (101,6.5);
  \draw[->] (79,11.5) -- (101,30.5);
  \draw[->] (79,18) -- (101,20);
  \draw[->] (79,33) -- (101,-6);
  \node[capt] at (90,-15) {\iflangja{(b) セグメンテーション}{(b) segmentation}};
\end{tikzpicture}
   (width ~118 mm)
\begin{tikzpicture}[x=1mm, y=1mm, font=\scriptsize\sffamily,
  >={Stealth[length=1.4mm]},
  pg/.style={draw, fill=ln-accent!22},
  sg/.style={draw, fill=ln-source!16},
  oth/.style={draw, fill=ln-box},
  lab/.style={font=\scriptsize\sffamily},
  hd/.style={font=\scriptsize\sffamily\bfseries, anchor=base},
  capt/.style={font=\footnotesize\sffamily}]
  % ================= (a) paging
  \node[hd] at (9,38) {\iflangja{仮想}{virtual}};
  \node[hd] at (45,38) {\iflangja{物理}{physical}};
  % virtual pages 0..5 (height 6)
  \foreach \j/\s in {0/pg, 1/pg, 2/pg, 4/pg, 5/pg} {
    \draw[\s] (2,6*\j) rectangle ++(14,6);
    \node[lab] at (9,3+6*\j) {\iflangja{ページ}{page} \j};
  }
  \draw[fill=white] (2,18) rectangle ++(14,6);
  \node[lab, text=ln-muted] at (9,21) {\iflangja{未使用}{unused}};
  % page frames 0..7 (height 4.5)
  \foreach \i in {0,...,7} { \draw[oth] (38,4.5*\i) rectangle ++(14,4.5); }
  \foreach \i in {0,2,5,7} { \draw[pg] (38,4.5*\i) rectangle ++(14,4.5); }
  \foreach \i in {1,3,4,6} { \node[lab, text=ln-muted] at (45,2.25+4.5*\i) {\iflangja{他}{other}}; }
  % disk
  \draw[fill=ln-box, rounded corners=2pt] (38,-11) rectangle (52,-4)
    node[midway, lab] {\iflangja{ディスク}{disk}};
  % mappings: page -> frame
  \foreach \j/\i in {0/5, 1/2, 4/7, 5/0} {
    \draw[->] (16,3+6*\j) -- (38,2.25+4.5*\i);
  }
  \draw[->] (16,15) -- (38,-7.5);
  \node[capt] at (27,-15) {\iflangja{(a) ページング}{(a) paging}};
  % ================= (b) segmentation
  \node[hd] at (72,38) {\iflangja{仮想}{virtual}};
  \node[hd] at (108,38) {\iflangja{物理}{physical}};
  \draw[sg] (65,0)  rectangle (79,9)  node[midway, lab] {\iflangja{コード}{code}};
  \draw[sg] (65,9)  rectangle (79,14) node[midway, lab] {\iflangja{データ}{data}};
  \draw[sg] (65,14) rectangle (79,22) node[midway, lab] {\iflangja{ヒープ}{heap}};
  \draw[fill=white] (65,22) rectangle (79,30);
  \draw[sg] (65,30) rectangle (79,36) node[midway, lab] {\iflangja{スタック}{stack}};
  % physical memory
  \draw[oth] (101,-11) rectangle (115,36);
  \draw[sg] (101,-9) rectangle (115,-3) node[midway, lab] {\iflangja{スタック}{stack}};
  \draw[sg] (101,2)  rectangle (115,11) node[midway, lab] {\iflangja{コード}{code}};
  \draw[sg] (101,16) rectangle (115,24) node[midway, lab] {\iflangja{ヒープ}{heap}};
  \draw[sg] (101,28) rectangle (115,33) node[midway, lab] {\iflangja{データ}{data}};
  \draw[->] (79,4.5) -- (101,6.5);
  \draw[->] (79,11.5) -- (101,30.5);
  \draw[->] (79,18) -- (101,20);
  \draw[->] (79,33) -- (101,-6);
  \node[capt] at (90,-15) {\iflangja{(b) セグメンテーション}{(b) segmentation}};
\end{tikzpicture}
   (width ~118 mm)
\begin{tikzpicture}[x=1mm, y=1mm, font=\scriptsize\sffamily,
  >={Stealth[length=1.4mm]},
  pg/.style={draw, fill=ln-accent!22},
  sg/.style={draw, fill=ln-source!16},
  oth/.style={draw, fill=ln-box},
  lab/.style={font=\scriptsize\sffamily},
  hd/.style={font=\scriptsize\sffamily\bfseries, anchor=base},
  capt/.style={font=\footnotesize\sffamily}]
  % ================= (a) paging
  \node[hd] at (9,38) {\iflangja{仮想}{virtual}};
  \node[hd] at (45,38) {\iflangja{物理}{physical}};
  % virtual pages 0..5 (height 6)
  \foreach \j/\s in {0/pg, 1/pg, 2/pg, 4/pg, 5/pg} {
    \draw[\s] (2,6*\j) rectangle ++(14,6);
    \node[lab] at (9,3+6*\j) {\iflangja{ページ}{page} \j};
  }
  \draw[fill=white] (2,18) rectangle ++(14,6);
  \node[lab, text=ln-muted] at (9,21) {\iflangja{未使用}{unused}};
  % page frames 0..7 (height 4.5)
  \foreach \i in {0,...,7} { \draw[oth] (38,4.5*\i) rectangle ++(14,4.5); }
  \foreach \i in {0,2,5,7} { \draw[pg] (38,4.5*\i) rectangle ++(14,4.5); }
  \foreach \i in {1,3,4,6} { \node[lab, text=ln-muted] at (45,2.25+4.5*\i) {\iflangja{他}{other}}; }
  % disk
  \draw[fill=ln-box, rounded corners=2pt] (38,-11) rectangle (52,-4)
    node[midway, lab] {\iflangja{ディスク}{disk}};
  % mappings: page -> frame
  \foreach \j/\i in {0/5, 1/2, 4/7, 5/0} {
    \draw[->] (16,3+6*\j) -- (38,2.25+4.5*\i);
  }
  \draw[->] (16,15) -- (38,-7.5);
  \node[capt] at (27,-15) {\iflangja{(a) ページング}{(a) paging}};
  % ================= (b) segmentation
  \node[hd] at (72,38) {\iflangja{仮想}{virtual}};
  \node[hd] at (108,38) {\iflangja{物理}{physical}};
  \draw[sg] (65,0)  rectangle (79,9)  node[midway, lab] {\iflangja{コード}{code}};
  \draw[sg] (65,9)  rectangle (79,14) node[midway, lab] {\iflangja{データ}{data}};
  \draw[sg] (65,14) rectangle (79,22) node[midway, lab] {\iflangja{ヒープ}{heap}};
  \draw[fill=white] (65,22) rectangle (79,30);
  \draw[sg] (65,30) rectangle (79,36) node[midway, lab] {\iflangja{スタック}{stack}};
  % physical memory
  \draw[oth] (101,-11) rectangle (115,36);
  \draw[sg] (101,-9) rectangle (115,-3) node[midway, lab] {\iflangja{スタック}{stack}};
  \draw[sg] (101,2)  rectangle (115,11) node[midway, lab] {\iflangja{コード}{code}};
  \draw[sg] (101,16) rectangle (115,24) node[midway, lab] {\iflangja{ヒープ}{heap}};
  \draw[sg] (101,28) rectangle (115,33) node[midway, lab] {\iflangja{データ}{data}};
  \draw[->] (79,4.5) -- (101,6.5);
  \draw[->] (79,11.5) -- (101,30.5);
  \draw[->] (79,18) -- (101,20);
  \draw[->] (79,33) -- (101,-6);
  \node[capt] at (90,-15) {\iflangja{(b) セグメンテーション}{(b) segmentation}};
\end{tikzpicture}

  \caption{(a) ページングは固定長のページをページフレームに対応付ける．
    ページ 2 はディスク上にあり，未使用のページには対応付けがない．
    (b) セグメンテーションは大きさの異なるセグメントを，同じ大きさの
    物理メモリの領域に対応付ける．}
  \label{fig:l08-goals}
\end{figure}

\section{ハードウェアによる支援}
\label{sec:l08-hardware}

すべてのメモリアクセスをソフトウェアで変換し検証するのは，あまりにも遅い．
\source{P3L2，メモリ管理：ハードウェアによる支援；OSTEP ch.~15，18--19．}
そこでメモリ管理は CPU パッケージ内のハードウェアに依存する．ハードウェアが
すべてのアクセスを処理し，オペレーティングシステムはデータ構造を用意して
例外を処理する（\cref{fig:l08-mmu}）．

\begin{definition}[メモリ管理ユニット]
  \term{MMU}[memory management unit]%
  \index{memory management unit|see{MMU}}（メモリ管理ユニット）とは，CPU が
  発行する仮想アドレスを物理アドレスへ変換し，変換が不可能なときにフォールト
  を起こす，CPU パッケージ内のハードウェア構成要素である．
\end{definition}

MMU がオペレーティングシステムにフォールトを報告するのは3つの場合である．
割り当てられていない仮想アドレスへのアクセス（不正アクセス），読み出し専用
メモリへの書き込みのように許可に違反するアクセス，そして物理メモリに存在
しておらず先にディスクから持ち込まなければならないページへのアクセスである．
MMU は，実行中のプロセスのページテーブルを指すレジスタを用いる．セグメン
テーションの場合には，セグメントのベースアドレス，リミット，その他の情報を
保持するレジスタを用いる．メモリ上のページテーブルを辿るのは遅いので，MMU は
有効な変換のキャッシュである TLB も備える（\cref{sec:l08-tlb}）．

変換そのものはすべてハードウェアで行われる．\caution{ハードウェアだけで変換
するアーキテクチャばかりではない．MIPS や SPARC では，TLB ミスが例外を起こ
し，カーネルのハンドラがソフトウェアで TLB を埋める．OSTEP ch.~19．}オペレー
ティングシステムは，
ハードウェアが使うページテーブルなどのデータ構造を作り，保守するが，その
形式はハードウェアが定め，したがってどのようなメモリ管理の機構が可能かも
ハードウェアが定める．例えば，ページの大きさやページテーブルエントリの
レイアウトである．

\begin{figure}[htbp]
  \centering
  % Hardware support for memory management: the MMU in the CPU package with
% its TLB and page table base register translates virtual addresses,
% walks the page tables in memory, and raises faults to the OS.
% Usage: % Hardware support for memory management: the MMU in the CPU package with
% its TLB and page table base register translates virtual addresses,
% walks the page tables in memory, and raises faults to the OS.
% Usage: % Hardware support for memory management: the MMU in the CPU package with
% its TLB and page table base register translates virtual addresses,
% walks the page tables in memory, and raises faults to the OS.
% Usage: \input{figures/l08-mmu}   (width ~116 mm)
\begin{tikzpicture}[x=1mm, y=1mm, font=\scriptsize\sffamily,
  >={Stealth[length=1.5mm]},
  box/.style={draw, fill=white, align=center, inner sep=1pt},
  hw/.style={draw, fill=ln-accent!14, align=center, inner sep=1pt},
  lab/.style={font=\scriptsize\sffamily, align=center},
  ttl/.style={font=\scriptsize\sffamily\bfseries, anchor=north west}]
  % CPU package
  \draw[ln-rule, dashed, rounded corners=3pt] (0,0) rectangle (64,36);
  \node[ttl, text=ln-muted] at (1,35.5) {\iflangja{CPU パッケージ}{CPU package}};
  \node[box, minimum width=14mm, minimum height=12mm] (cpu) at (10,17) {CPU};
  \draw[fill=ln-box] (30,3) rectangle (61,31);
  \node[ttl] at (30.5,30.5) {MMU};
  \node[hw, minimum width=27mm, minimum height=7mm] (tlb) at (45.5,19)
    {\iflangja{TLB（変換のキャッシュ）}{TLB (translation cache)}};
  \node[hw, minimum width=27mm, minimum height=7mm] (reg) at (45.5,9)
    {\iflangja{ページテーブル\\ベースレジスタ}{page table base\\register}};
  \draw[->] (cpu.east) -- (30,17) node[pos=0.5, above, lab] {\iflangja{仮想\\アドレス}{virtual\\address}};
  % physical memory
  \draw[fill=ln-box] (86,0) rectangle (116,36);
  \node[ttl] at (86.5,35.5) {\iflangja{物理メモリ}{physical memory}};
  \node[hw, minimum width=24mm, minimum height=7mm] (pt) at (101,24)
    {\iflangja{ページテーブル}{page tables}};
  \node[box, minimum width=24mm, minimum height=11mm] (pages) at (101,8)
    {\iflangja{プロセスの\\ページ}{pages of\\processes}};
  \draw[->] (61,8) -- (pages.west) node[pos=0.5, above, lab] {\iflangja{物理アドレス}{physical address}};
  \draw[<->, densely dashed] (61,24) -- (pt.west) node[pos=0.5, above, lab] {\iflangja{ページテーブル参照}{page table walk}};
  % faults
  \node[box, minimum width=31mm, minimum height=7mm, fill=ln-caution!8] (os) at (45.5,-12)
    {\iflangja{OS のフォールトハンドラ}{OS fault handler}};
  \draw[->, ln-caution] (45.5,3) -- (os.north) node[pos=0.45, right, lab, text=ln-caution]
    {\iflangja{フォールト}{fault}};
\end{tikzpicture}
   (width ~116 mm)
\begin{tikzpicture}[x=1mm, y=1mm, font=\scriptsize\sffamily,
  >={Stealth[length=1.5mm]},
  box/.style={draw, fill=white, align=center, inner sep=1pt},
  hw/.style={draw, fill=ln-accent!14, align=center, inner sep=1pt},
  lab/.style={font=\scriptsize\sffamily, align=center},
  ttl/.style={font=\scriptsize\sffamily\bfseries, anchor=north west}]
  % CPU package
  \draw[ln-rule, dashed, rounded corners=3pt] (0,0) rectangle (64,36);
  \node[ttl, text=ln-muted] at (1,35.5) {\iflangja{CPU パッケージ}{CPU package}};
  \node[box, minimum width=14mm, minimum height=12mm] (cpu) at (10,17) {CPU};
  \draw[fill=ln-box] (30,3) rectangle (61,31);
  \node[ttl] at (30.5,30.5) {MMU};
  \node[hw, minimum width=27mm, minimum height=7mm] (tlb) at (45.5,19)
    {\iflangja{TLB（変換のキャッシュ）}{TLB (translation cache)}};
  \node[hw, minimum width=27mm, minimum height=7mm] (reg) at (45.5,9)
    {\iflangja{ページテーブル\\ベースレジスタ}{page table base\\register}};
  \draw[->] (cpu.east) -- (30,17) node[pos=0.5, above, lab] {\iflangja{仮想\\アドレス}{virtual\\address}};
  % physical memory
  \draw[fill=ln-box] (86,0) rectangle (116,36);
  \node[ttl] at (86.5,35.5) {\iflangja{物理メモリ}{physical memory}};
  \node[hw, minimum width=24mm, minimum height=7mm] (pt) at (101,24)
    {\iflangja{ページテーブル}{page tables}};
  \node[box, minimum width=24mm, minimum height=11mm] (pages) at (101,8)
    {\iflangja{プロセスの\\ページ}{pages of\\processes}};
  \draw[->] (61,8) -- (pages.west) node[pos=0.5, above, lab] {\iflangja{物理アドレス}{physical address}};
  \draw[<->, densely dashed] (61,24) -- (pt.west) node[pos=0.5, above, lab] {\iflangja{ページテーブル参照}{page table walk}};
  % faults
  \node[box, minimum width=31mm, minimum height=7mm, fill=ln-caution!8] (os) at (45.5,-12)
    {\iflangja{OS のフォールトハンドラ}{OS fault handler}};
  \draw[->, ln-caution] (45.5,3) -- (os.north) node[pos=0.45, right, lab, text=ln-caution]
    {\iflangja{フォールト}{fault}};
\end{tikzpicture}
   (width ~116 mm)
\begin{tikzpicture}[x=1mm, y=1mm, font=\scriptsize\sffamily,
  >={Stealth[length=1.5mm]},
  box/.style={draw, fill=white, align=center, inner sep=1pt},
  hw/.style={draw, fill=ln-accent!14, align=center, inner sep=1pt},
  lab/.style={font=\scriptsize\sffamily, align=center},
  ttl/.style={font=\scriptsize\sffamily\bfseries, anchor=north west}]
  % CPU package
  \draw[ln-rule, dashed, rounded corners=3pt] (0,0) rectangle (64,36);
  \node[ttl, text=ln-muted] at (1,35.5) {\iflangja{CPU パッケージ}{CPU package}};
  \node[box, minimum width=14mm, minimum height=12mm] (cpu) at (10,17) {CPU};
  \draw[fill=ln-box] (30,3) rectangle (61,31);
  \node[ttl] at (30.5,30.5) {MMU};
  \node[hw, minimum width=27mm, minimum height=7mm] (tlb) at (45.5,19)
    {\iflangja{TLB（変換のキャッシュ）}{TLB (translation cache)}};
  \node[hw, minimum width=27mm, minimum height=7mm] (reg) at (45.5,9)
    {\iflangja{ページテーブル\\ベースレジスタ}{page table base\\register}};
  \draw[->] (cpu.east) -- (30,17) node[pos=0.5, above, lab] {\iflangja{仮想\\アドレス}{virtual\\address}};
  % physical memory
  \draw[fill=ln-box] (86,0) rectangle (116,36);
  \node[ttl] at (86.5,35.5) {\iflangja{物理メモリ}{physical memory}};
  \node[hw, minimum width=24mm, minimum height=7mm] (pt) at (101,24)
    {\iflangja{ページテーブル}{page tables}};
  \node[box, minimum width=24mm, minimum height=11mm] (pages) at (101,8)
    {\iflangja{プロセスの\\ページ}{pages of\\processes}};
  \draw[->] (61,8) -- (pages.west) node[pos=0.5, above, lab] {\iflangja{物理アドレス}{physical address}};
  \draw[<->, densely dashed] (61,24) -- (pt.west) node[pos=0.5, above, lab] {\iflangja{ページテーブル参照}{page table walk}};
  % faults
  \node[box, minimum width=31mm, minimum height=7mm, fill=ln-caution!8] (os) at (45.5,-12)
    {\iflangja{OS のフォールトハンドラ}{OS fault handler}};
  \draw[->, ln-caution] (45.5,3) -- (os.north) node[pos=0.45, right, lab, text=ln-caution]
    {\iflangja{フォールト}{fault}};
\end{tikzpicture}

  \caption{メモリ管理のためのハードウェアによる支援．MMU は TLB の助けを
    借りて仮想アドレスを変換し，変換がキャッシュされていないときは物理メモリ
    上のページテーブルを辿り，アクセスが不可能なときはオペレーティング
    システムへフォールトを起こす．}
  \label{fig:l08-mmu}
\end{figure}

\section{ページテーブル}
\label{sec:l08-page-tables}

ページテーブル（第2章）は，仮想アドレス空間のすべてのページを，それを保持
するページフレームに対応付ける．\source{P3L2，ページテーブル；OSTEP
ch.~18；Silberschatz ら ch.~9．}\index{ぺーじてーぶる@ページテーブル}ページと
ページフレームは同じ大きさなので，ページテーブルはアドレスを1つずつ対応付け
る必要はない．ページ内でのアドレスの位置は仮想メモリでも物理メモリでも同じ
であり，変換が必要なのはページだけだからである．

$2^{d}$ バイトのページでは，仮想アドレスは2つの部分に分けられる．
\begin{description}
  \item[仮想ページ番号] 上位のビットが
    \term[かそうぺーじばんごう]{仮想ページ番号}[virtual page number]%
    \index{VPN|see{仮想ページ番号}}（VPN）を成し，これがページテーブルの
    索引となる．
  \item[オフセット] 下位 $d$ ビットが\term[おふせっと]{オフセット}[offset]を
    成す．これはページ内でのアドレスの位置であり，変換されない．
\end{description}
ページテーブルの各エントリを
\term[ぺーじてーぶるえんとり]{ページテーブルエントリ}[page table entry]%
\index{PTE|see{ページテーブルエントリ}}（PTE）と呼ぶ．PTE は，他の情報
（\cref{sec:l08-pte}）とともに，そのページを含むフレームの
\term[ぶつりふれーむばんごう]{物理フレーム番号}[physical frame number]%
\index{PFN|see{物理フレーム番号}}（PFN）を保持する．\caution{PFN の幅は VPN
の幅と同じとは限らない．物理アドレス拡張（PAE）を用いる 32 ビットの x86
は，20 ビットの VPN を 24 ビットの PFN へ対応付け，\SI{64}{\gibi\byte} の物
理メモリを扱う．}物理アドレスは，PFN と
変換されないオフセットを組み合わせて得られる（\cref{fig:l08-translation}）．
\begin{equation}
  \text{VPN} = \left\lfloor \frac{\mathit{VA}}{2^{d}} \right\rfloor ,
  \qquad
  \text{offset} = \mathit{VA} \bmod 2^{d} ,
  \qquad
  \mathit{PA} = \text{PFN} \cdot 2^{d} + \text{offset} .
  \label{eq:l08-translation}
\end{equation}

\begin{figure}[htbp]
  \centering
  % Address translation with a (flat) page table: the VPN indexes the page
% table found through the page table base register; the PFN of the entry
% is combined with the unchanged offset.
% Usage: % Address translation with a (flat) page table: the VPN indexes the page
% table found through the page table base register; the PFN of the entry
% is combined with the unchanged offset.
% Usage: % Address translation with a (flat) page table: the VPN indexes the page
% table found through the page table base register; the PFN of the entry
% is combined with the unchanged offset.
% Usage: \input{figures/l08-translation}   (width ~116 mm)
\begin{tikzpicture}[x=1mm, y=1mm, font=\scriptsize\sffamily,
  >={Stealth[length=1.5mm]},
  cell/.style={draw, minimum height=7mm, inner sep=0pt, anchor=south west},
  lab/.style={font=\scriptsize\sffamily, align=center},
  hd/.style={font=\scriptsize\sffamily\bfseries, anchor=south}]
  % virtual address
  \node[hd, anchor=south west] at (0,37.5) {\iflangja{仮想アドレス}{virtual address}};
  \node[cell, minimum width=18mm, fill=ln-accent!22] (vpn) at (0,30) {VPN};
  \node[cell, minimum width=16mm, fill=ln-box] (voff) at (18,30) {\iflangja{オフセット}{offset}};
  % page table base register
  \node[cell, minimum width=24mm, minimum height=6mm, fill=ln-source!14, align=center]
    (ptbr) at (46,36) {\iflangja{ページテーブル\\ベースレジスタ}{page table\\base register}};
  % page table
  \foreach \k in {0,...,5} {
    \draw[fill=white] (46,5*\k) rectangle ++(16,5);
    \draw[fill=white] (62,5*\k) rectangle ++(8,5);
  }
  \draw[fill=ln-accent!22] (46,15) rectangle ++(16,5) node[midway] {PFN};
  \draw[fill=ln-accent!10] (62,15) rectangle ++(8,5) node[midway, font=\iflangja{\tiny}{\scriptsize}\sffamily] {\iflangja{フラグ}{flags}};
  \node[lab, anchor=north] at (58,-0.5) {\iflangja{ページテーブル}{page table}};
  \draw[->] (58,36) -- (58,30);
  % VPN selects the entry
  \draw[->] (9,30) |- (46,17.5) node[pos=0.75, above, lab] {\iflangja{索引}{index}};
  % physical address
  \node[hd, anchor=north] at (99,13.5) {\iflangja{物理アドレス}{physical address}};
  \node[cell, minimum width=18mm, fill=ln-accent!22] (pfn) at (82,14) {PFN};
  \node[cell, minimum width=16mm, fill=ln-box] (poff) at (100,14) {\iflangja{オフセット}{offset}};
  \draw[->] (70,17.5) -- (82,17.5);
  % offset copied unchanged
  \draw[->, densely dashed] (26,37) -- (26,49) -| (108,21)
    node[pos=0.25, above, lab] {\iflangja{オフセットはそのまま}{offset copied unchanged}};
\end{tikzpicture}
   (width ~116 mm)
\begin{tikzpicture}[x=1mm, y=1mm, font=\scriptsize\sffamily,
  >={Stealth[length=1.5mm]},
  cell/.style={draw, minimum height=7mm, inner sep=0pt, anchor=south west},
  lab/.style={font=\scriptsize\sffamily, align=center},
  hd/.style={font=\scriptsize\sffamily\bfseries, anchor=south}]
  % virtual address
  \node[hd, anchor=south west] at (0,37.5) {\iflangja{仮想アドレス}{virtual address}};
  \node[cell, minimum width=18mm, fill=ln-accent!22] (vpn) at (0,30) {VPN};
  \node[cell, minimum width=16mm, fill=ln-box] (voff) at (18,30) {\iflangja{オフセット}{offset}};
  % page table base register
  \node[cell, minimum width=24mm, minimum height=6mm, fill=ln-source!14, align=center]
    (ptbr) at (46,36) {\iflangja{ページテーブル\\ベースレジスタ}{page table\\base register}};
  % page table
  \foreach \k in {0,...,5} {
    \draw[fill=white] (46,5*\k) rectangle ++(16,5);
    \draw[fill=white] (62,5*\k) rectangle ++(8,5);
  }
  \draw[fill=ln-accent!22] (46,15) rectangle ++(16,5) node[midway] {PFN};
  \draw[fill=ln-accent!10] (62,15) rectangle ++(8,5) node[midway, font=\iflangja{\tiny}{\scriptsize}\sffamily] {\iflangja{フラグ}{flags}};
  \node[lab, anchor=north] at (58,-0.5) {\iflangja{ページテーブル}{page table}};
  \draw[->] (58,36) -- (58,30);
  % VPN selects the entry
  \draw[->] (9,30) |- (46,17.5) node[pos=0.75, above, lab] {\iflangja{索引}{index}};
  % physical address
  \node[hd, anchor=north] at (99,13.5) {\iflangja{物理アドレス}{physical address}};
  \node[cell, minimum width=18mm, fill=ln-accent!22] (pfn) at (82,14) {PFN};
  \node[cell, minimum width=16mm, fill=ln-box] (poff) at (100,14) {\iflangja{オフセット}{offset}};
  \draw[->] (70,17.5) -- (82,17.5);
  % offset copied unchanged
  \draw[->, densely dashed] (26,37) -- (26,49) -| (108,21)
    node[pos=0.25, above, lab] {\iflangja{オフセットはそのまま}{offset copied unchanged}};
\end{tikzpicture}
   (width ~116 mm)
\begin{tikzpicture}[x=1mm, y=1mm, font=\scriptsize\sffamily,
  >={Stealth[length=1.5mm]},
  cell/.style={draw, minimum height=7mm, inner sep=0pt, anchor=south west},
  lab/.style={font=\scriptsize\sffamily, align=center},
  hd/.style={font=\scriptsize\sffamily\bfseries, anchor=south}]
  % virtual address
  \node[hd, anchor=south west] at (0,37.5) {\iflangja{仮想アドレス}{virtual address}};
  \node[cell, minimum width=18mm, fill=ln-accent!22] (vpn) at (0,30) {VPN};
  \node[cell, minimum width=16mm, fill=ln-box] (voff) at (18,30) {\iflangja{オフセット}{offset}};
  % page table base register
  \node[cell, minimum width=24mm, minimum height=6mm, fill=ln-source!14, align=center]
    (ptbr) at (46,36) {\iflangja{ページテーブル\\ベースレジスタ}{page table\\base register}};
  % page table
  \foreach \k in {0,...,5} {
    \draw[fill=white] (46,5*\k) rectangle ++(16,5);
    \draw[fill=white] (62,5*\k) rectangle ++(8,5);
  }
  \draw[fill=ln-accent!22] (46,15) rectangle ++(16,5) node[midway] {PFN};
  \draw[fill=ln-accent!10] (62,15) rectangle ++(8,5) node[midway, font=\iflangja{\tiny}{\scriptsize}\sffamily] {\iflangja{フラグ}{flags}};
  \node[lab, anchor=north] at (58,-0.5) {\iflangja{ページテーブル}{page table}};
  \draw[->] (58,36) -- (58,30);
  % VPN selects the entry
  \draw[->] (9,30) |- (46,17.5) node[pos=0.75, above, lab] {\iflangja{索引}{index}};
  % physical address
  \node[hd, anchor=north] at (99,13.5) {\iflangja{物理アドレス}{physical address}};
  \node[cell, minimum width=18mm, fill=ln-accent!22] (pfn) at (82,14) {PFN};
  \node[cell, minimum width=16mm, fill=ln-box] (poff) at (100,14) {\iflangja{オフセット}{offset}};
  \draw[->] (70,17.5) -- (82,17.5);
  % offset copied unchanged
  \draw[->, densely dashed] (26,37) -- (26,49) -| (108,21)
    node[pos=0.25, above, lab] {\iflangja{オフセットはそのまま}{offset copied unchanged}};
\end{tikzpicture}

  \caption{ページテーブルによるアドレス変換．VPN が，実行中のプロセスの
    ページテーブルのエントリを選ぶ．そのエントリの PFN と，変換されない
    オフセットとを合わせたものが物理アドレスになる．}
  \label{fig:l08-translation}
\end{figure}

\begin{example}
  32 ビットの仮想アドレスと \SI{4}{\kibi\byte} のページ（$d = 12$）では，
  VPN は 20 ビット，オフセットは 12 ビット，すなわち 16 進で 3 桁である．
  仮想アドレス \texttt{0x00035A7C} の VPN は \texttt{0x00035}，オフセットは
  \texttt{0xA7C} である．ページテーブルのエントリ \texttt{0x35} が PFN
  \texttt{0x001C2} を保持していれば，物理アドレスは \texttt{0x001C2A7C} と
  なる．\texttt{0x00035000} から \texttt{0x00035FFF} までのすべてのアドレスが
  同じエントリで変換される．
\end{example}

\subsection{初回アクセス時の割り当てと回収}

プロセスが，例えば大きな配列のために仮想アドレスの範囲を予約しても，その
時点では物理メモリは割り当てられない．割り当てられるのは，プロセスがその
メモリにアクセスしたときである．この方針を
\term[しょかいあくせすじのわりあて]{初回アクセス時の割り当て}[allocation on first touch]と
呼ぶ．ページへの最初のアクセスはフォールトを起こし，そこでオペレーティング
システムがそのページにページフレームを割り当て，ページテーブルに対応付けを
作る．

\needspace{8\baselineskip}
\begin{example}
  次のコードは \SI{256}{\mebi\byte} を予約するが，ページフレームを受け取るの
  は，ループが $2^{28}/2^{12} = 65\,536$ 個のページに1つずつアクセスして
  いくのに応じてである．
  \begin{lstlisting}[language=C, numbers=none]
size_t n = 256UL << 20;          /* 256 MiB */
char *buf = malloc(n);           /* まだフレームなし */
for (size_t i = 0; i < n; i += 4096)
    buf[i] = 1;                  /* 初回アクセスでフォールト */
\end{lstlisting}
\end{example}

逆に，あるページが長く使われていないとき，オペレーティングシステムはその
フレームを回収してよい．ページの内容をディスクへ移し，対応付けを無効と記す
（\cref{sec:l08-pte}）．プロセスがそのページに再びアクセスすると，MMU は
フォールトを起こしてオペレーティングシステムへトラップする．オペレーティング
システムは，そのアクセスが許可されていることを確かめ，内容を，場合によっては
別のページフレームへ，メモリに戻し，対応付けを作り直す．プロセスは，遅延の
ほかには何も気づかないまま実行を続ける．

オペレーティングシステムはプロセスごとに1つのページテーブルを保持する．MMU
は，実行中のプロセスのページテーブルを
\term[ぺーじてーぶるべーすれじすた]{ページテーブルベースレジスタ}[page table base register]を
通じて見つける．これはそのページテーブルのアドレスを保持する CPU のレジスタ
であり，x86 ではレジスタ CR3 である．コンテキストスイッチの際，オペレー
ティングシステムは次のプロセスのページテーブルのアドレスをこのレジスタに
読み込む．

\section{ページテーブルエントリ}
\label{sec:l08-pte}

PFN のほかに，\source{P3L2，ページテーブルエントリ；OSTEP ch.~18，21；
Silberschatz ら ch.~9．}どの PTE もフラグを含む．MMU はアクセスのたびに
これを検査し，オペレーティングシステムはこれを使ってメモリを管理する．
最も重要なフラグは次のものである．
\begin{description}
  \item[有効ビット] \term[ゆうこうびっと]{有効ビット}[valid bit]は，
    \emph{存在ビット}とも呼ばれ，対応付けが有効かどうか，すなわちエントリが
    指すフレームにそのページが物理メモリ上に存在するかどうかを示す．
  \item[ダーティビット] ハードウェアは，ページが書き込まれるたびに
    \term[だーてぃびっと]{ダーティビット}[dirty bit]を立てる．例えばファイルの
    内容をキャッシュするページは，ダーティビットが立っているときにだけ
    ファイルへ書き戻せばよい．
  \item[アクセスビット] ハードウェアは，ページが読み出しでも書き込みでも
    アクセスされるたびに\term[あくせすびっと]{アクセスビット}[access bit]
    （参照ビット）を立てる．
  \item[保護ビット] \term[ほごびっと]{保護ビット}[protection bits]は，その
    ページに対して許される操作---読み出し，書き込み，実行---を指定する．
\end{description}
オペレーティングシステムはダーティビットとアクセスビットを読み，それらに
基づいて処置を行った後にクリアする．

32 ビットの x86 プロセッサでは PTE は 32 ビットである（\cref{fig:l08-pte}）．
上位 20 ビットが PFN を保持し，下位のビットには存在ビット P，ダーティビット
D，アクセスビット A が含まれる．R/W は保護ビットであり，0 ならページは
読み出し専用，1 なら書き込みもできる．U/S は特権レベルを制限し，0 なら
ページはスーパバイザ（カーネル）モードでのみ，1 ならユーザモードでも
アクセスできる．PWT と PCD はキャッシュに関する情報（ライトスルー，
キャッシュ禁止）を伝え，ビット 9--11 はハードウェアが使わず，オペレー
ティングシステムが自由に使える．

\begin{figure}[htbp]
  \centering
  % Layout of a 32-bit x86 page table entry (4 KiB page, no PAE).
% Usage: % Layout of a 32-bit x86 page table entry (4 KiB page, no PAE).
% Usage: % Layout of a 32-bit x86 page table entry (4 KiB page, no PAE).
% Usage: \input{figures/l08-pte}   (width ~116 mm)
\begin{tikzpicture}[x=1mm, y=1mm, font=\scriptsize\sffamily,
  num/.style={font=\tiny\sffamily, text=ln-muted, anchor=south},
  grp/.style={font=\scriptsize\sffamily, anchor=north, align=center}]
  % field boundaries: PFN 0..34, unused 34..46, then 9 single bits of 7.8 mm
  \draw[fill=ln-accent!22] (0,0) rectangle (34,7) node[midway] {\iflangja{物理フレーム番号}{physical frame number}};
  \draw[fill=white] (34,0) rectangle (46,7) node[midway] {\iflangja{未使用}{unused}};
  \foreach \k/\n/\f in {0/G/white, 1/PAT/white, 2/D/ln-accent!22, 3/A/ln-accent!22,
                        4/PCD/ln-box, 5/PWT/ln-box, 6/{U/S}/ln-source!16,
                        7/{R/W}/ln-source!16, 8/P/ln-accent!22} {
    \draw[fill=\f] (46+7.8*\k,0) rectangle ++(7.8,7);
    \node at (49.9+7.8*\k,3.5) {\n};
  }
  % bit numbers
  \node[num] at (1.5,7.2) {31};
  \node[num] at (32.5,7.2) {12};
  \node[num] at (35.5,7.2) {11};
  \node[num] at (44.5,7.2) {9};
  \foreach \k/\b in {0/8, 1/7, 2/6, 3/5, 4/4, 5/3, 6/2, 7/1, 8/0} {
    \node[num] at (49.9+7.8*\k,7.2) {\b};
  }
  % groups
  \draw[ln-muted] (62.2,-1) -- (62.2,-2) -- (76.6,-2) -- (76.6,-1);
  \node[grp] at (69.4,-2.5) {\iflangja{ダーティ，\\アクセス}{dirty,\\accessed}};
  \draw[ln-muted] (77.8,-1) -- (77.8,-2) -- (92.2,-2) -- (92.2,-1);
  \node[grp] at (85,-2.5) {\iflangja{キャッシュ\\関連}{caching}};
  \draw[ln-muted] (93.4,-1) -- (93.4,-2) -- (107.8,-2) -- (107.8,-1);
  \node[grp] at (100.6,-2.5) {\iflangja{保護}{protection}};
  \draw[ln-muted] (112.3,-1) -- (112.3,-2);
  \node[grp] at (112.3,-2.5) {\iflangja{存在}{present}};
\end{tikzpicture}
   (width ~116 mm)
\begin{tikzpicture}[x=1mm, y=1mm, font=\scriptsize\sffamily,
  num/.style={font=\tiny\sffamily, text=ln-muted, anchor=south},
  grp/.style={font=\scriptsize\sffamily, anchor=north, align=center}]
  % field boundaries: PFN 0..34, unused 34..46, then 9 single bits of 7.8 mm
  \draw[fill=ln-accent!22] (0,0) rectangle (34,7) node[midway] {\iflangja{物理フレーム番号}{physical frame number}};
  \draw[fill=white] (34,0) rectangle (46,7) node[midway] {\iflangja{未使用}{unused}};
  \foreach \k/\n/\f in {0/G/white, 1/PAT/white, 2/D/ln-accent!22, 3/A/ln-accent!22,
                        4/PCD/ln-box, 5/PWT/ln-box, 6/{U/S}/ln-source!16,
                        7/{R/W}/ln-source!16, 8/P/ln-accent!22} {
    \draw[fill=\f] (46+7.8*\k,0) rectangle ++(7.8,7);
    \node at (49.9+7.8*\k,3.5) {\n};
  }
  % bit numbers
  \node[num] at (1.5,7.2) {31};
  \node[num] at (32.5,7.2) {12};
  \node[num] at (35.5,7.2) {11};
  \node[num] at (44.5,7.2) {9};
  \foreach \k/\b in {0/8, 1/7, 2/6, 3/5, 4/4, 5/3, 6/2, 7/1, 8/0} {
    \node[num] at (49.9+7.8*\k,7.2) {\b};
  }
  % groups
  \draw[ln-muted] (62.2,-1) -- (62.2,-2) -- (76.6,-2) -- (76.6,-1);
  \node[grp] at (69.4,-2.5) {\iflangja{ダーティ，\\アクセス}{dirty,\\accessed}};
  \draw[ln-muted] (77.8,-1) -- (77.8,-2) -- (92.2,-2) -- (92.2,-1);
  \node[grp] at (85,-2.5) {\iflangja{キャッシュ\\関連}{caching}};
  \draw[ln-muted] (93.4,-1) -- (93.4,-2) -- (107.8,-2) -- (107.8,-1);
  \node[grp] at (100.6,-2.5) {\iflangja{保護}{protection}};
  \draw[ln-muted] (112.3,-1) -- (112.3,-2);
  \node[grp] at (112.3,-2.5) {\iflangja{存在}{present}};
\end{tikzpicture}
   (width ~116 mm)
\begin{tikzpicture}[x=1mm, y=1mm, font=\scriptsize\sffamily,
  num/.style={font=\tiny\sffamily, text=ln-muted, anchor=south},
  grp/.style={font=\scriptsize\sffamily, anchor=north, align=center}]
  % field boundaries: PFN 0..34, unused 34..46, then 9 single bits of 7.8 mm
  \draw[fill=ln-accent!22] (0,0) rectangle (34,7) node[midway] {\iflangja{物理フレーム番号}{physical frame number}};
  \draw[fill=white] (34,0) rectangle (46,7) node[midway] {\iflangja{未使用}{unused}};
  \foreach \k/\n/\f in {0/G/white, 1/PAT/white, 2/D/ln-accent!22, 3/A/ln-accent!22,
                        4/PCD/ln-box, 5/PWT/ln-box, 6/{U/S}/ln-source!16,
                        7/{R/W}/ln-source!16, 8/P/ln-accent!22} {
    \draw[fill=\f] (46+7.8*\k,0) rectangle ++(7.8,7);
    \node at (49.9+7.8*\k,3.5) {\n};
  }
  % bit numbers
  \node[num] at (1.5,7.2) {31};
  \node[num] at (32.5,7.2) {12};
  \node[num] at (35.5,7.2) {11};
  \node[num] at (44.5,7.2) {9};
  \foreach \k/\b in {0/8, 1/7, 2/6, 3/5, 4/4, 5/3, 6/2, 7/1, 8/0} {
    \node[num] at (49.9+7.8*\k,7.2) {\b};
  }
  % groups
  \draw[ln-muted] (62.2,-1) -- (62.2,-2) -- (76.6,-2) -- (76.6,-1);
  \node[grp] at (69.4,-2.5) {\iflangja{ダーティ，\\アクセス}{dirty,\\accessed}};
  \draw[ln-muted] (77.8,-1) -- (77.8,-2) -- (92.2,-2) -- (92.2,-1);
  \node[grp] at (85,-2.5) {\iflangja{キャッシュ\\関連}{caching}};
  \draw[ln-muted] (93.4,-1) -- (93.4,-2) -- (107.8,-2) -- (107.8,-1);
  \node[grp] at (100.6,-2.5) {\iflangja{保護}{protection}};
  \draw[ln-muted] (112.3,-1) -- (112.3,-2);
  \node[grp] at (112.3,-2.5) {\iflangja{存在}{present}};
\end{tikzpicture}

  \caption{32 ビット x86 のページテーブルエントリ．グローバルビット G と
    ページ属性ビット PAT はここでは扱わない．}
  \label{fig:l08-pte}
\end{figure}

\subsection{ページフォールト}
\label{sec:l08-page-fault}

MMU は PTE を，アドレスの変換にもアクセスの検査にも使う．

\begin{definition}[ページフォールト]
  \term[ぺーじふぉーると]{ページフォールト}[page fault]とは，現在のページ
  テーブルエントリではアクセスを実行できないときに MMU が起こす例外である．
  エントリが有効でないか，アクセスがその保護ビットに違反する場合である．
\end{definition}

ページフォールトが起こると，CPU はエラーコードをカーネルスタックに置いて
カーネルへトラップし，カーネルはページフォールトハンドラを実行する．x86 では
エラーコードが，ページが存在しなかったのかアクセスが保護ビットに違反したの
か，アクセスが読み出しだったのか書き込みだったのか，ユーザモードとカーネル
モードのどちらから来たのかを伝え，フォールトした仮想アドレスはレジスタ CR2
に保存される．ハンドラは，この情報と，自分が保持するアドレス空間の記録とから，
取るべき動作を決める（\cref{tab:l08-faults}）．

\begin{table}[htbp]
  \centering
  \caption{ページテーブルエントリとオペレーティングシステム自身の記録に
    応じた，アクセスの扱い．}
  \label{tab:l08-faults}
  \small
  \begin{tabular}{@{}lll@{}}
    \toprule
    PTE & 状況 & 動作 \\
    \midrule
    有効   & アクセスは許可         & なし：MMU が変換する \\
    無効   & 割り当て済み，未アクセス & フレームを割り当てる（初回アクセス） \\
    無効   & ページはディスク上     & 持ち込む（\cref{sec:l08-demand}） \\
    無効   & アドレスは未割り当て   & 不正アクセス：\texttt{SIGSEGV} \\
    有効   & アクセスが保護に違反   & \texttt{SIGSEGV}，またはコピー（\cref{sec:l08-cow}） \\
    \bottomrule
  \end{tabular}
\end{table}

不正なアクセスを行ったプロセスはシグナル \texttt{SIGSEGV}（第5章）を受け
取り，その既定の動作はプロセスを終了させる．他のフォールトはすべてオペレー
ティングシステムが解決し，その後フォールトした命令が再び実行される．

\section{ページテーブルの大きさ}
\label{sec:l08-pt-size}

平坦なページテーブルは，ページが使われているかどうかにかかわらず，すべての
仮想ページ番号に対して1つのエントリを持つ．\source{P3L2，ページテーブルの
大きさ；OSTEP ch.~20．}$n$ ビットの仮想アドレス，$2^{d}$ バイトのページ，
$e$ バイトのページテーブルエントリでは，エントリ数とテーブルの大きさは
\begin{equation}
  N_{\mathrm{PTE}} = \frac{2^{n}}{2^{d}} = 2^{\,n-d} ,
  \qquad
  S_{\mathrm{PT}} = 2^{\,n-d} \cdot e
  \label{eq:l08-pt-size}
\end{equation}
である．\SI{4}{\kibi\byte} のページと 4 バイトのエントリを持つ 32 ビット
アーキテクチャでは，ページテーブルは $2^{20}$ 個のエントリを持ち
\SI{4}{\mebi\byte} を占める．64 ビットのアドレスと 8 バイトのエントリでは，
1つのプロセスにつき $2^{52}$ 個のエントリ，$2^{55}$ バイト，すなわち
\SI{32}{\pebi\byte} を占めることになる．どのプロセスも自分のページテーブルを
持つので，この代価にはさらにプロセスの数が掛かる．しかもプロセスはアドレス空間の
ごく一部しか使わないので，これらのエントリのほとんどは，決して使われない
ページを記述することになる．\tip{ページが大きいほどエントリは少ない．ページ
の大きさを2倍にすれば $N_{\mathrm{PTE}}$ は半分になる
（\cref{sec:l08-page-size}）．}

\begin{example}
  x86-64 アーキテクチャは 48 ビットの仮想アドレスを使う．
  \SI{4}{\kibi\byte} のページと 8 バイトのエントリでは，平坦なページテーブル
  はプロセスごとに $2^{36}$ 個のエントリ，$2^{39}$ バイト，すなわち
  \SI{512}{\gibi\byte} を必要とする．
\end{example}

\section{多段ページテーブル}
\label{sec:l08-multilevel}

したがって，大きなアドレス空間に対して平坦なページテーブルは実用的でない．
\source{P3L2，多段ページテーブル；OSTEP ch.~20；Silberschatz ら ch.~9．}
アドレス空間の大部分は使われていないので，ページテーブルは代わりに階層的に
構成され，大きな未使用領域のページにはページテーブルエントリが不要になる．

\begin{definition}[多段ページテーブル]
  \term[ただんぺーじてーぶる]{多段ページテーブル}[multi-level page table]とは，
  テーブルの木である．外側のテーブルである
  \term[ぺーじでぃれくとり]{ページディレクトリ}[page directory]は内側の
  ページテーブルへのポインタを保持し，内側のページテーブルが PFN を保持する．
  内側のページテーブルは，有効なページを含むアドレス空間の領域に対してだけ
  割り当てられる．
\end{definition}

2段のテーブルでは，仮想アドレスは3つの部分に分けられる
（\cref{fig:l08-multilevel}）．$p_1$ がページディレクトリの索引，$p_2$ が
ディレクトリのエントリが指す内側のページテーブルの索引となり，オフセット
$d$ がフレーム内のアドレスを選ぶ．$p_2$ が $k$ ビットなら，内側のページ
テーブル1つが仮想アドレス空間の $2^{k+d}$ バイトを覆う．\tip{8 バイトの
エントリ $2^{9}$ 個のテーブルはちょうど \SI{4}{\kibi\byte} である．1段
当たり 9 ビットの索引にすると，どの段のテーブルも1つのページフレームに
収まる．}この大きさの
アドレス空間の隙間には内側のページテーブルはまったく不要であり，その
ディレクトリエントリを無効と記すだけでよい．プロセスが後にその領域にメモリを
割り当てたら，オペレーティングシステムはそのための内側のページテーブルを
割り当て，ディレクトリエントリを埋める．

\begin{figure}[htbp]
  \centering
  % Two-level page table: p1 indexes the page directory, whose entry points
% to an inner page table; p2 indexes that table, whose entry gives the page
% frame; the offset d selects the byte within the frame.  Directory entries
% without an inner table are left white.
% Usage: % Two-level page table: p1 indexes the page directory, whose entry points
% to an inner page table; p2 indexes that table, whose entry gives the page
% frame; the offset d selects the byte within the frame.  Directory entries
% without an inner table are left white.
% Usage: % Two-level page table: p1 indexes the page directory, whose entry points
% to an inner page table; p2 indexes that table, whose entry gives the page
% frame; the offset d selects the byte within the frame.  Directory entries
% without an inner table are left white.
% Usage: \input{figures/l08-multilevel}   (width ~118 mm)
\begin{tikzpicture}[x=1mm, y=1mm, font=\scriptsize\sffamily,
  >={Stealth[length=1.5mm]},
  cell/.style={draw, minimum height=6mm, inner sep=0pt, anchor=south west},
  lab/.style={font=\scriptsize\sffamily, align=center},
  idx/.style={font=\scriptsize\sffamily, text=ln-muted},
  hd/.style={font=\scriptsize\sffamily\bfseries, anchor=south, align=center}]
  % page table base register
  \node[cell, minimum width=22mm, minimum height=7mm, fill=ln-source!14, align=center]
    at (0,38) {\iflangja{ページテーブル\\ベースレジスタ}{page table\\base register}};
  \draw[->] (22,41.5) -| (39,36);
  % page directory: rows top to bottom, centres 33, 27, 21, 15, 9
  \node[hd] at (39,44) {\iflangja{ページ\\ディレクトリ}{page\\directory}};
  \foreach \r/\f in {0/ln-accent!10, 1/ln-accent!30, 2/ln-accent!10, 3/white, 4/white} {
    \draw[fill=\f] (30,30-6*\r) rectangle ++(18,6);
  }
  \node[lab] at (39,27) {\iflangja{内部表へ}{$\to$ inner}};
  % inner page table
  \node[hd] at (73,44) {\iflangja{内部\\ページテーブル}{inner\\page table}};
  \foreach \r/\f in {0/ln-accent!10, 1/white, 2/ln-accent!10, 3/ln-accent!30, 4/white} {
    \draw[fill=\f] (64,30-6*\r) rectangle ++(18,6);
  }
  \node[lab] at (73,15) {PFN};
  % page frame
  \node[hd] at (107,44) {\iflangja{ページ\\フレーム}{page\\frame}};
  \draw[fill=ln-box] (98,6) rectangle (116,36);
  \draw[fill=ln-accent!30] (98,17) rectangle (116,20);
  % pointers along the top
  \draw[->] (48,27) -- (56,27) |- (73,40.5) -- (73,36);
  \draw[->] (82,15) -- (90,15) |- (107,40.5) -- (107,36);
  % virtual address
  \node[cell, minimum width=16mm, fill=ln-accent!22] (p1) at (31,-16) {$p_1$};
  \node[cell, minimum width=16mm, fill=ln-accent!22] (p2) at (47,-16) {$p_2$};
  \node[cell, minimum width=16mm, fill=ln-box] (d) at (63,-16) {$d$};
  \node[lab, anchor=east] at (30.5,-13) {\iflangja{仮想アドレス}{virtual\\address}};
  % indices from below
  \draw[->] (39,-10) -- (39,6);
  \draw[->] (55,-10) |- (73,-2) -- (73,6);
  \draw[->] (71,-10) |- (107,-6) -- (107,17);
  \node[idx, anchor=west] at (39.5,1) {\iflangja{索引}{index}};
  \node[idx, anchor=west] at (73.5,2) {\iflangja{索引}{index}};
  \node[idx, anchor=west] at (107.5,2) {\iflangja{オフセット}{offset}};
\end{tikzpicture}
   (width ~118 mm)
\begin{tikzpicture}[x=1mm, y=1mm, font=\scriptsize\sffamily,
  >={Stealth[length=1.5mm]},
  cell/.style={draw, minimum height=6mm, inner sep=0pt, anchor=south west},
  lab/.style={font=\scriptsize\sffamily, align=center},
  idx/.style={font=\scriptsize\sffamily, text=ln-muted},
  hd/.style={font=\scriptsize\sffamily\bfseries, anchor=south, align=center}]
  % page table base register
  \node[cell, minimum width=22mm, minimum height=7mm, fill=ln-source!14, align=center]
    at (0,38) {\iflangja{ページテーブル\\ベースレジスタ}{page table\\base register}};
  \draw[->] (22,41.5) -| (39,36);
  % page directory: rows top to bottom, centres 33, 27, 21, 15, 9
  \node[hd] at (39,44) {\iflangja{ページ\\ディレクトリ}{page\\directory}};
  \foreach \r/\f in {0/ln-accent!10, 1/ln-accent!30, 2/ln-accent!10, 3/white, 4/white} {
    \draw[fill=\f] (30,30-6*\r) rectangle ++(18,6);
  }
  \node[lab] at (39,27) {\iflangja{内部表へ}{$\to$ inner}};
  % inner page table
  \node[hd] at (73,44) {\iflangja{内部\\ページテーブル}{inner\\page table}};
  \foreach \r/\f in {0/ln-accent!10, 1/white, 2/ln-accent!10, 3/ln-accent!30, 4/white} {
    \draw[fill=\f] (64,30-6*\r) rectangle ++(18,6);
  }
  \node[lab] at (73,15) {PFN};
  % page frame
  \node[hd] at (107,44) {\iflangja{ページ\\フレーム}{page\\frame}};
  \draw[fill=ln-box] (98,6) rectangle (116,36);
  \draw[fill=ln-accent!30] (98,17) rectangle (116,20);
  % pointers along the top
  \draw[->] (48,27) -- (56,27) |- (73,40.5) -- (73,36);
  \draw[->] (82,15) -- (90,15) |- (107,40.5) -- (107,36);
  % virtual address
  \node[cell, minimum width=16mm, fill=ln-accent!22] (p1) at (31,-16) {$p_1$};
  \node[cell, minimum width=16mm, fill=ln-accent!22] (p2) at (47,-16) {$p_2$};
  \node[cell, minimum width=16mm, fill=ln-box] (d) at (63,-16) {$d$};
  \node[lab, anchor=east] at (30.5,-13) {\iflangja{仮想アドレス}{virtual\\address}};
  % indices from below
  \draw[->] (39,-10) -- (39,6);
  \draw[->] (55,-10) |- (73,-2) -- (73,6);
  \draw[->] (71,-10) |- (107,-6) -- (107,17);
  \node[idx, anchor=west] at (39.5,1) {\iflangja{索引}{index}};
  \node[idx, anchor=west] at (73.5,2) {\iflangja{索引}{index}};
  \node[idx, anchor=west] at (107.5,2) {\iflangja{オフセット}{offset}};
\end{tikzpicture}
   (width ~118 mm)
\begin{tikzpicture}[x=1mm, y=1mm, font=\scriptsize\sffamily,
  >={Stealth[length=1.5mm]},
  cell/.style={draw, minimum height=6mm, inner sep=0pt, anchor=south west},
  lab/.style={font=\scriptsize\sffamily, align=center},
  idx/.style={font=\scriptsize\sffamily, text=ln-muted},
  hd/.style={font=\scriptsize\sffamily\bfseries, anchor=south, align=center}]
  % page table base register
  \node[cell, minimum width=22mm, minimum height=7mm, fill=ln-source!14, align=center]
    at (0,38) {\iflangja{ページテーブル\\ベースレジスタ}{page table\\base register}};
  \draw[->] (22,41.5) -| (39,36);
  % page directory: rows top to bottom, centres 33, 27, 21, 15, 9
  \node[hd] at (39,44) {\iflangja{ページ\\ディレクトリ}{page\\directory}};
  \foreach \r/\f in {0/ln-accent!10, 1/ln-accent!30, 2/ln-accent!10, 3/white, 4/white} {
    \draw[fill=\f] (30,30-6*\r) rectangle ++(18,6);
  }
  \node[lab] at (39,27) {\iflangja{内部表へ}{$\to$ inner}};
  % inner page table
  \node[hd] at (73,44) {\iflangja{内部\\ページテーブル}{inner\\page table}};
  \foreach \r/\f in {0/ln-accent!10, 1/white, 2/ln-accent!10, 3/ln-accent!30, 4/white} {
    \draw[fill=\f] (64,30-6*\r) rectangle ++(18,6);
  }
  \node[lab] at (73,15) {PFN};
  % page frame
  \node[hd] at (107,44) {\iflangja{ページ\\フレーム}{page\\frame}};
  \draw[fill=ln-box] (98,6) rectangle (116,36);
  \draw[fill=ln-accent!30] (98,17) rectangle (116,20);
  % pointers along the top
  \draw[->] (48,27) -- (56,27) |- (73,40.5) -- (73,36);
  \draw[->] (82,15) -- (90,15) |- (107,40.5) -- (107,36);
  % virtual address
  \node[cell, minimum width=16mm, fill=ln-accent!22] (p1) at (31,-16) {$p_1$};
  \node[cell, minimum width=16mm, fill=ln-accent!22] (p2) at (47,-16) {$p_2$};
  \node[cell, minimum width=16mm, fill=ln-box] (d) at (63,-16) {$d$};
  \node[lab, anchor=east] at (30.5,-13) {\iflangja{仮想アドレス}{virtual\\address}};
  % indices from below
  \draw[->] (39,-10) -- (39,6);
  \draw[->] (55,-10) |- (73,-2) -- (73,6);
  \draw[->] (71,-10) |- (107,-6) -- (107,17);
  \node[idx, anchor=west] at (39.5,1) {\iflangja{索引}{index}};
  \node[idx, anchor=west] at (73.5,2) {\iflangja{索引}{index}};
  \node[idx, anchor=west] at (107.5,2) {\iflangja{オフセット}{offset}};
\end{tikzpicture}

  \caption{2段のページテーブル．白いディレクトリエントリには内側のページ
    テーブルがない．変換には，ページディレクトリへの1回と内側のページ
    テーブルへの1回のアクセスを要する．}
  \label{fig:l08-multilevel}
\end{figure}

\begin{example}
  \label{ex:l08-multilevel}
  20 ビットの仮想アドレスと \SI{1}{\kibi\byte} のページ（$d = 10$）を考える．
  平坦なページテーブルは $2^{10} = 1024$ 個のエントリを持つ．$p_1$ が 4
  ビット，$p_2$ が 6 ビットの2段のテーブルでは，ページディレクトリは 16 個の
  エントリを持ち，64 個のエントリからなる内側のテーブルはそれぞれ
  \SI{64}{\kibi\byte} を覆う．あるプロセスが最下位の \SI{100}{\kibi\byte}
  （テキスト，データ，ヒープ，ページ 0--99）と最上位の \SI{20}{\kibi\byte}
  （スタック，ページ 1004--1023）を使うとする．このプロセスには，ページ
  0--63，64--127，960--1023 のための内側のテーブルが必要で，合計
  $16 + 3 \cdot 64 = 208$ 個のエントリ，平坦なテーブルの約5分の1である．
\end{example}

階層はさらに段を加えて拡張できる．3段目はページディレクトリへのポインタの
テーブルを，4段目はそのテーブルへのポインタのテーブルを加える．\margin{x86-64
は 12 ビットのオフセットの上に 9 ビットずつ4段を置く．$4\cdot 9 + 12 = 48$
ビットのアドレスである．}段を増やすことが最も効くのは 64 ビットアーキテク
チャであり，そのアドレス空間はより大きく，かつより疎である．木の上位にある
1つのエントリが覆う領域が大きいほど，その下にテーブルを必要としない隙間も
大きくなる．

代価は変換のレイテンシである．どの段のエントリも順にメモリから読まれるので，
1段ごとにアドレスの変換に1回のメモリアクセスが加わる．段が多いほどページ
テーブルのためのメモリは節約できるが，1回の変換は遅くなる．\margin{x86-64 が
実装するのは 64 ビットのうち 48 ビットだけである．Intel の5段ページングは5段
目を加えてこれを 57 ビット（\SI{128}{\pebi\byte}）まで広げる．Linux は 4.14
以降これを支援する．}

\section{変換の高速化：TLB}
\label{sec:l08-tlb}

1段のページテーブルでは，どのメモリ参照もページテーブルエントリへの1回と
データへの1回の計2回のメモリアクセスを要し，4段では5回を要する．
\source{P3L2，変換の速度：TLB；OSTEP ch.~19；Silberschatz ら ch.~9．}MMU は
これらのアクセスの大半をキャッシュで避ける．

\begin{definition}[変換索引バッファ]
  \term{TLB}[translation lookaside buffer]%
  \index{translation lookaside buffer|see{TLB}}（変換索引バッファ）とは，MMU に
  組み込まれた，最近使われたアドレス変換のキャッシュである．TLB ヒットの際は
  ページテーブルにアクセスせずにアドレスが変換され，TLB ミスの際はページ
  テーブルが辿られ，その変換が TLB に格納される．
\end{definition}

TLB のエントリは，その変換の保護ビットと有効ビットも保持するので，ページ
テーブルなしにアクセスを検査できる．\margin{MMU は上位の段のエントリもキャッ
シュするので（paging-structure caches），ミスの代価が $L$ 回のアクセスに達す
ることはまれである．Intel SDM vol.~3A，4.10.3．}メモリアクセスには\emph{局所
性}があるので，
少数のエントリで高いヒット率が得られる．今アクセスされたページは間もなく
再びアクセスされる可能性が高く（時間的局所性），それに近いアドレス---たいて
い同じページの中にある---が次にアクセスされる可能性が高い（空間的局所性）．
例えば Sandy Bridge 世代の Intel Core i7 プロセッサは，コアごとに 64 エントリ
のデータ TLB，128 エントリの命令 TLB，および両方の変換が共有する 512 エントリ
の第2段 TLB を持つ（いずれも \SI{4}{\kibi\byte} のページに対して）．
\margin{コンテキストスイッチは TLB のエントリを古くする．アドレス空間で
タグ付けされていない限り，エントリはフラッシュされる（第12章）．}

TLB の参照時間を $t_{\mathrm{TLB}}$，メモリアクセス時間を
$t_{\mathrm{mem}}$，ヒット率を $h$，ページテーブルを $L$ 段とすると，メモリ
アクセスの実効時間は
\begin{equation}
  t_{\mathrm{eff}} = h\,(t_{\mathrm{TLB}} + t_{\mathrm{mem}})
    + (1-h)\,\bigl(t_{\mathrm{TLB}} + (L+1)\,t_{\mathrm{mem}}\bigr)
  \label{eq:l08-eat}
\end{equation}
である．

\begin{example}
  $t_{\mathrm{mem}} = \SI{80}{\nano\second}$，
  $t_{\mathrm{TLB}} = \SI{1}{\nano\second}$，$L = 4$ とす
  る．\tip{\cref{eq:l08-eat}は $t_{\mathrm{eff}} = t_{\mathrm{TLB}} +
  \bigl(1 + L(1-h)\bigr)\,t_{\mathrm{mem}}$ と書き直せる．
  参照の代価は常に払い，テーブルを辿る代価はミスのときだけ払う．}TLB がなけ
  れば，
  どのアクセスも $5 \cdot 80 = \SI{400}{\nano\second}$ かかる．ヒット率が
  \SI{99}{\percent} なら，\cref{eq:l08-eat} より $0.99 \cdot 81 + 0.01 \cdot
  401 \approx \SI{84.2}{\nano\second}$ となり，変換のないアクセスより
  \SI{5}{\percent} 増えるだけである．ヒット率が \SI{90}{\percent} なら
  $0.9 \cdot 81 + 0.1 \cdot 401 = \SI{113}{\nano\second}$ となる．
\end{example}

\begin{keypoint}
  平坦なページテーブルは仮想アドレス空間に比例し，大きすぎる．多段ページ
  テーブルは使用中の領域に対してのみテーブルを割り当て，その代価は1段当たり
  1回のメモリアクセスである．局所性を利用する TLB が，ほとんどすべての変換で
  そのアクセスを省く．
\end{keypoint}

\section{逆ページテーブル}
\label{sec:l08-inverted}

ページテーブルを小さく保つもう1つの方法は，仮想メモリではなく物理メモリに
従ってテーブルを構成することである．\source{P3L2，逆ページテーブル；OSTEP
ch.~20；Silberschatz ら ch.~9．}物理メモリは，すべてのプロセスの仮想アドレス
空間を合わせたものよりはるかに小さく，どのページフレームも同時にはたかだか
1つのページしか保持しない．

\begin{definition}[逆ページテーブル]
  \term[ぎゃくぺーじてーぶる]{逆ページテーブル}[inverted page table]は，
  物理メモリのページフレームごとに1つのエントリを持つ．フレーム $i$ の
  エントリは，フレーム $i$ が保持するページのプロセス識別子と仮想ページ番号を
  記録する．
\end{definition}

システム全体で逆ページテーブルは1つで済み，その大きさは物理メモリに比例し，
プロセスの数にもそのアドレス空間の大きさにも依存しない．\caution{フレームご
とに1つのエントリということは，フレームごとに $(\text{pid}, p)$ の組が1つと
いうことである．複数のプロセスが共有するフレームはこの構造に収まらず，
共有メモリの扱いが難しくなる．}仮想アドレスは，
プロセス識別子（pid），ページ番号 $p$，オフセット $d$ からなる．これを変換
するには，$(\text{pid}, p)$ に一致するエントリをテーブルから探す．その
エントリの索引 $i$ がフレーム番号であり，物理アドレスは $i$ と $d$ から
作られる（\cref{fig:l08-inverted}\,(a)）．

\begin{figure}[htbp]
  \centering
  % (a) Inverted page table: one entry per page frame; the entry matching
% (pid, p) is searched, and its index i is the frame number.
% (b) Hashed page table: a hash of (pid, p) selects a bucket whose chain
% holds (pid, p, frame) entries.
% Usage: % (a) Inverted page table: one entry per page frame; the entry matching
% (pid, p) is searched, and its index i is the frame number.
% (b) Hashed page table: a hash of (pid, p) selects a bucket whose chain
% holds (pid, p, frame) entries.
% Usage: % (a) Inverted page table: one entry per page frame; the entry matching
% (pid, p) is searched, and its index i is the frame number.
% (b) Hashed page table: a hash of (pid, p) selects a bucket whose chain
% holds (pid, p, frame) entries.
% Usage: \input{figures/l08-inverted}   (width ~118 mm)
\begin{tikzpicture}[x=1mm, y=1mm, font=\scriptsize\sffamily,
  >={Stealth[length=1.5mm]},
  cell/.style={draw, minimum height=6mm, inner sep=0pt, anchor=south west},
  lab/.style={font=\scriptsize\sffamily, align=center},
  idx/.style={font=\scriptsize\sffamily, text=ln-muted},
  hd/.style={font=\scriptsize\sffamily\bfseries, anchor=south, align=center},
  capt/.style={font=\footnotesize\sffamily, anchor=west}]
  % ================= (a) inverted page table
  \node[capt] at (0,42) {\iflangja{(a) 逆ページテーブル}{(a) inverted page table}};
  \node[lab, anchor=south west] at (0,27) {\iflangja{仮想アドレス}{virtual address}};
  \node[cell, minimum width=10mm, fill=ln-source!16] at (0,20) {pid};
  \node[cell, minimum width=10mm, fill=ln-accent!22] at (10,20) {$p$};
  \node[cell, minimum width=10mm, fill=ln-box] at (20,20) {$d$};
  \node[hd] at (58,31) {\iflangja{フレームごとに 1 エントリ}{one entry per frame}};
  \foreach \r/\a/\b in {0/7/3, 1/2/12, 2/5/9, 4/7/8} {
    \draw[fill=white] (48,24-6*\r) rectangle ++(10,6) node[midway] {\a};
    \draw[fill=white] (58,24-6*\r) rectangle ++(10,6) node[midway] {\b};
  }
  \draw[fill=ln-accent!22] (48,6) rectangle ++(10,6) node[midway] {pid};
  \draw[fill=ln-accent!22] (58,6) rectangle ++(10,6) node[midway] {$p$};
  \node[idx, anchor=east] at (47.5,9) {$i$};
  \node[idx] at (58,-1.5) {\iflangja{（索引 $i$ ＝ フレーム番号）}{(index $i$ = frame number)}};
  \draw (5,20) -- (5,14) -- (15,14);
  \draw[->] (15,20) |- (44,9) node[pos=0.75, above, lab] {\iflangja{探索}{search}};
  % physical address
  \node[cell, minimum width=10mm, fill=ln-accent!22] at (86,6) {$i$};
  \node[cell, minimum width=10mm, fill=ln-box] at (96,6) {$d$};
  \node[lab, anchor=north west] at (86,5.5) {\iflangja{物理アドレス}{physical address}};
  \draw[->] (68,9) -- (86,9);
  \draw[->, densely dashed] (25,26) -- (25,37) -| (101,12);
  % ================= (b) hashed page table
  \node[capt] at (0,-8) {\iflangja{(b) ハッシュページテーブル}{(b) hashed page table}};
  \node[cell, minimum width=10mm, fill=ln-source!16] at (0,-24) {pid};
  \node[cell, minimum width=10mm, fill=ln-accent!22] at (10,-24) {$p$};
  \node[cell, minimum width=10mm, fill=ln-box] at (20,-24) {$d$};
  \node[draw, circle, inner sep=1pt, fill=white, font=\scriptsize\sffamily] (h) at (34,-21) {$h$};
  \draw (5,-18) -- (5,-14) -- (34,-14);
  \draw (15,-18) -- (15,-14);
  \draw[->] (34,-14) -- (h.north);
  % hash table buckets
  \foreach \r in {0,...,4} { \draw[fill=white] (44,-12-5*\r) rectangle ++(8,5); }
  \draw[fill=ln-accent!22] (44,-27) rectangle ++(8,5);
  \node[idx] at (48,-34.5) {\iflangja{バケット}{buckets}};
  \draw[->] (h.east) -- (44,-24.5);
  % chain entries: pid | p | frame
  \foreach \x/\a/\b/\c in {60/4/31/17, 88/2/31/6} {
    \draw[fill=white] (\x,-27) rectangle ++(6,5) node[midway] {\a};
    \draw[fill=white] (\x+6,-27) rectangle ++(6,5) node[midway] {\b};
    \draw[fill=ln-accent!10] (\x+12,-27) rectangle ++(6,5) node[midway] {\c};
    \draw[fill=ln-box] (\x+18,-27) rectangle ++(3,5);
  }
  \draw[->] (48,-24.5) -- (60,-24.5);
  \fill (79.5,-24.5) circle (0.5);
  \draw[->] (79.5,-24.5) -- (88,-24.5);
  \node[idx, anchor=north] at (69,-27.5) {pid\quad $p$\quad \iflangja{枠}{frame}};
  \node[idx, anchor=north] at (97,-27.5) {pid\quad $p$\quad \iflangja{枠}{frame}};
  \node[lab, text width=40mm, anchor=north west] at (60,-33)
    {\iflangja{連鎖を (pid, $p$) で照合}{chain searched for (pid, $p$)}};
\end{tikzpicture}
   (width ~118 mm)
\begin{tikzpicture}[x=1mm, y=1mm, font=\scriptsize\sffamily,
  >={Stealth[length=1.5mm]},
  cell/.style={draw, minimum height=6mm, inner sep=0pt, anchor=south west},
  lab/.style={font=\scriptsize\sffamily, align=center},
  idx/.style={font=\scriptsize\sffamily, text=ln-muted},
  hd/.style={font=\scriptsize\sffamily\bfseries, anchor=south, align=center},
  capt/.style={font=\footnotesize\sffamily, anchor=west}]
  % ================= (a) inverted page table
  \node[capt] at (0,42) {\iflangja{(a) 逆ページテーブル}{(a) inverted page table}};
  \node[lab, anchor=south west] at (0,27) {\iflangja{仮想アドレス}{virtual address}};
  \node[cell, minimum width=10mm, fill=ln-source!16] at (0,20) {pid};
  \node[cell, minimum width=10mm, fill=ln-accent!22] at (10,20) {$p$};
  \node[cell, minimum width=10mm, fill=ln-box] at (20,20) {$d$};
  \node[hd] at (58,31) {\iflangja{フレームごとに 1 エントリ}{one entry per frame}};
  \foreach \r/\a/\b in {0/7/3, 1/2/12, 2/5/9, 4/7/8} {
    \draw[fill=white] (48,24-6*\r) rectangle ++(10,6) node[midway] {\a};
    \draw[fill=white] (58,24-6*\r) rectangle ++(10,6) node[midway] {\b};
  }
  \draw[fill=ln-accent!22] (48,6) rectangle ++(10,6) node[midway] {pid};
  \draw[fill=ln-accent!22] (58,6) rectangle ++(10,6) node[midway] {$p$};
  \node[idx, anchor=east] at (47.5,9) {$i$};
  \node[idx] at (58,-1.5) {\iflangja{（索引 $i$ ＝ フレーム番号）}{(index $i$ = frame number)}};
  \draw (5,20) -- (5,14) -- (15,14);
  \draw[->] (15,20) |- (44,9) node[pos=0.75, above, lab] {\iflangja{探索}{search}};
  % physical address
  \node[cell, minimum width=10mm, fill=ln-accent!22] at (86,6) {$i$};
  \node[cell, minimum width=10mm, fill=ln-box] at (96,6) {$d$};
  \node[lab, anchor=north west] at (86,5.5) {\iflangja{物理アドレス}{physical address}};
  \draw[->] (68,9) -- (86,9);
  \draw[->, densely dashed] (25,26) -- (25,37) -| (101,12);
  % ================= (b) hashed page table
  \node[capt] at (0,-8) {\iflangja{(b) ハッシュページテーブル}{(b) hashed page table}};
  \node[cell, minimum width=10mm, fill=ln-source!16] at (0,-24) {pid};
  \node[cell, minimum width=10mm, fill=ln-accent!22] at (10,-24) {$p$};
  \node[cell, minimum width=10mm, fill=ln-box] at (20,-24) {$d$};
  \node[draw, circle, inner sep=1pt, fill=white, font=\scriptsize\sffamily] (h) at (34,-21) {$h$};
  \draw (5,-18) -- (5,-14) -- (34,-14);
  \draw (15,-18) -- (15,-14);
  \draw[->] (34,-14) -- (h.north);
  % hash table buckets
  \foreach \r in {0,...,4} { \draw[fill=white] (44,-12-5*\r) rectangle ++(8,5); }
  \draw[fill=ln-accent!22] (44,-27) rectangle ++(8,5);
  \node[idx] at (48,-34.5) {\iflangja{バケット}{buckets}};
  \draw[->] (h.east) -- (44,-24.5);
  % chain entries: pid | p | frame
  \foreach \x/\a/\b/\c in {60/4/31/17, 88/2/31/6} {
    \draw[fill=white] (\x,-27) rectangle ++(6,5) node[midway] {\a};
    \draw[fill=white] (\x+6,-27) rectangle ++(6,5) node[midway] {\b};
    \draw[fill=ln-accent!10] (\x+12,-27) rectangle ++(6,5) node[midway] {\c};
    \draw[fill=ln-box] (\x+18,-27) rectangle ++(3,5);
  }
  \draw[->] (48,-24.5) -- (60,-24.5);
  \fill (79.5,-24.5) circle (0.5);
  \draw[->] (79.5,-24.5) -- (88,-24.5);
  \node[idx, anchor=north] at (69,-27.5) {pid\quad $p$\quad \iflangja{枠}{frame}};
  \node[idx, anchor=north] at (97,-27.5) {pid\quad $p$\quad \iflangja{枠}{frame}};
  \node[lab, text width=40mm, anchor=north west] at (60,-33)
    {\iflangja{連鎖を (pid, $p$) で照合}{chain searched for (pid, $p$)}};
\end{tikzpicture}
   (width ~118 mm)
\begin{tikzpicture}[x=1mm, y=1mm, font=\scriptsize\sffamily,
  >={Stealth[length=1.5mm]},
  cell/.style={draw, minimum height=6mm, inner sep=0pt, anchor=south west},
  lab/.style={font=\scriptsize\sffamily, align=center},
  idx/.style={font=\scriptsize\sffamily, text=ln-muted},
  hd/.style={font=\scriptsize\sffamily\bfseries, anchor=south, align=center},
  capt/.style={font=\footnotesize\sffamily, anchor=west}]
  % ================= (a) inverted page table
  \node[capt] at (0,42) {\iflangja{(a) 逆ページテーブル}{(a) inverted page table}};
  \node[lab, anchor=south west] at (0,27) {\iflangja{仮想アドレス}{virtual address}};
  \node[cell, minimum width=10mm, fill=ln-source!16] at (0,20) {pid};
  \node[cell, minimum width=10mm, fill=ln-accent!22] at (10,20) {$p$};
  \node[cell, minimum width=10mm, fill=ln-box] at (20,20) {$d$};
  \node[hd] at (58,31) {\iflangja{フレームごとに 1 エントリ}{one entry per frame}};
  \foreach \r/\a/\b in {0/7/3, 1/2/12, 2/5/9, 4/7/8} {
    \draw[fill=white] (48,24-6*\r) rectangle ++(10,6) node[midway] {\a};
    \draw[fill=white] (58,24-6*\r) rectangle ++(10,6) node[midway] {\b};
  }
  \draw[fill=ln-accent!22] (48,6) rectangle ++(10,6) node[midway] {pid};
  \draw[fill=ln-accent!22] (58,6) rectangle ++(10,6) node[midway] {$p$};
  \node[idx, anchor=east] at (47.5,9) {$i$};
  \node[idx] at (58,-1.5) {\iflangja{（索引 $i$ ＝ フレーム番号）}{(index $i$ = frame number)}};
  \draw (5,20) -- (5,14) -- (15,14);
  \draw[->] (15,20) |- (44,9) node[pos=0.75, above, lab] {\iflangja{探索}{search}};
  % physical address
  \node[cell, minimum width=10mm, fill=ln-accent!22] at (86,6) {$i$};
  \node[cell, minimum width=10mm, fill=ln-box] at (96,6) {$d$};
  \node[lab, anchor=north west] at (86,5.5) {\iflangja{物理アドレス}{physical address}};
  \draw[->] (68,9) -- (86,9);
  \draw[->, densely dashed] (25,26) -- (25,37) -| (101,12);
  % ================= (b) hashed page table
  \node[capt] at (0,-8) {\iflangja{(b) ハッシュページテーブル}{(b) hashed page table}};
  \node[cell, minimum width=10mm, fill=ln-source!16] at (0,-24) {pid};
  \node[cell, minimum width=10mm, fill=ln-accent!22] at (10,-24) {$p$};
  \node[cell, minimum width=10mm, fill=ln-box] at (20,-24) {$d$};
  \node[draw, circle, inner sep=1pt, fill=white, font=\scriptsize\sffamily] (h) at (34,-21) {$h$};
  \draw (5,-18) -- (5,-14) -- (34,-14);
  \draw (15,-18) -- (15,-14);
  \draw[->] (34,-14) -- (h.north);
  % hash table buckets
  \foreach \r in {0,...,4} { \draw[fill=white] (44,-12-5*\r) rectangle ++(8,5); }
  \draw[fill=ln-accent!22] (44,-27) rectangle ++(8,5);
  \node[idx] at (48,-34.5) {\iflangja{バケット}{buckets}};
  \draw[->] (h.east) -- (44,-24.5);
  % chain entries: pid | p | frame
  \foreach \x/\a/\b/\c in {60/4/31/17, 88/2/31/6} {
    \draw[fill=white] (\x,-27) rectangle ++(6,5) node[midway] {\a};
    \draw[fill=white] (\x+6,-27) rectangle ++(6,5) node[midway] {\b};
    \draw[fill=ln-accent!10] (\x+12,-27) rectangle ++(6,5) node[midway] {\c};
    \draw[fill=ln-box] (\x+18,-27) rectangle ++(3,5);
  }
  \draw[->] (48,-24.5) -- (60,-24.5);
  \fill (79.5,-24.5) circle (0.5);
  \draw[->] (79.5,-24.5) -- (88,-24.5);
  \node[idx, anchor=north] at (69,-27.5) {pid\quad $p$\quad \iflangja{枠}{frame}};
  \node[idx, anchor=north] at (97,-27.5) {pid\quad $p$\quad \iflangja{枠}{frame}};
  \node[lab, text width=40mm, anchor=north west] at (60,-33)
    {\iflangja{連鎖を (pid, $p$) で照合}{chain searched for (pid, $p$)}};
\end{tikzpicture}

  \caption{(a) 逆ページテーブルを $(\text{pid}, p)$ で探索する．一致した
    エントリの索引がフレーム番号である．(b) ハッシュページテーブルは，
    ハッシュ関数 $h$ が選んだバケットの連鎖だけを探索する．}
  \label{fig:l08-inverted}
\end{figure}

テーブルの線形探索は遅い．変換の大半は TLB が処理するが，ミスしたときの探索も
やはり速くなければならず，これはハッシュ法で加速される．
\term[はっしゅぺーじてーぶる]{ハッシュページテーブル}[hashed page table]では，
ページ番号にプロセス識別子を合わせてハッシュ関数を適用し，バケットを選ぶ．
バケットは，キーがそのバケットにハッシュされるエントリの連結リストの先頭で
ある（\cref{fig:l08-inverted}\,(b)）．各要素はプロセス識別子，ページ番号，
フレーム番号，および次の要素へのポインタを保持するので，1回の変換は1本の
短いリストだけを探索する．\margin{ハッシュ化した逆ページテーブルは IBM
AS/400，64 ビットの PowerPC，Itanium で使われた．x86 は一貫して多段ページ
テーブルである．}

\begin{example}
  \SI{16}{\gibi\byte} の物理メモリと \SI{4}{\kibi\byte} のフレームを持つ
  マシンは $2^{34}/2^{12} = 2^{22}$ 個のフレームを持つ．16 バイトのエントリ
  なら，その逆ページテーブルは $2^{26}$ バイト，すなわち \SI{64}{\mebi\byte}
  を占める．マシンが 10 個のプロセスを動かしていようと 1000 個を動かして
  いようと同じである．
\end{example}

\section{セグメンテーション}
\label{sec:l08-segmentation}

セグメンテーションでは，アドレス空間は任意の大きさのセグメントに分割され，
各セグメントは，コード，データ，ヒープ，スタックのような，プロセスの論理的に
意味のある構成要素に対応する．\source{P3L2，セグメンテーション；OSTEP
ch.~16；Silberschatz ら ch.~9；Bovet と Cesati ch.~2．}仮想アドレスは，x86
では\emph{論理アドレス}と呼ばれ，\emph{セグメントセレクタ}とセグメント内の
オフセットからなる．最も単純な形では，セグメントは物理メモリの連続した領域で
あり，それが始まるアドレスである\emph{ベース}と，その大きさである
\emph{リミット}によって記述される．物理アドレスはベースにオフセットを加えた
ものであり，オフセットがリミット未満のときにだけ有効である．そうでなければ
アクセスはフォールトする．

\begin{example}
  あるプロセスが，ベース \texttt{0x4000}，リミット \texttt{0x1800} のコード
  セグメントと，ベース \texttt{0x9000}，リミット \texttt{0x0800} のスタック
  セグメントを持つとする．論理アドレス（コード，\texttt{0x0A20}）は
  $\texttt{0x0A20} < \texttt{0x1800}$ なので \texttt{0x4A20} に変換される．
  論理アドレス（スタック，\texttt{0x0900}）は，そのオフセットがスタック
  セグメントのリミットを超えるのでフォールトする．
\end{example}

x86 では，セレクタはセグメントレジスタに保持される．コード用の CS，データ用の
DS，スタック用の SS などである．セレクタは\emph{セグメント記述子表}---グ
ローバル記述子表（GDT）またはローカル記述子表（LDT）---の索引であり，その
エントリであるセグメント記述子が，セグメントのベース，リミット，および許可を
保持する．\margin{記述子は 8 バイトである．セレクタは 13 ビットの索引---ゆえ
に 8192 エントリ---と，GDT か LDT かを選ぶビット，2 ビットの特権レベルを持
つ．}ベースとオフセットの和が\emph{リニアアドレス}である．セグメン
テーションはページングと組み合わされる．リニアアドレスはページテーブルに
よって物理アドレスへ変換される（\cref{fig:l08-segmentation}）．
\begin{itemize}
  \item 32 ビットの x86（IA-32）では両方の機構が有効である．記述子表は
    最大 8192 個のセグメント記述子を保持する．Linux は，多くのオペレー
    ティングシステムと同様に，アドレス 0 から始まりアドレス空間全体にわたる
    セグメントを定義し，変換と保護はページングに任せる．
  \item x86-64 では，セグメンテーションは後方互換のためにのみ残されている．
    64 ビットモードでは，プロセッサはセグメントのベースを 0 として扱い
    （FS と GS を除く），リミットも検査しないので，アドレスの変換はページング
    だけが行う．\sidenote{2つの例外は実際に使われている．Linux は，
    実行中のスレッドのスレッドローカル記憶を FS のベースに，
    カーネルの CPU ごとのデータを GS のベースに置く．システムコールでカーネ
    ルに入るとき，命令 \texttt{swapgs} が後者のユーザ側とカーネル側の値を入
    れ替える．}
\end{itemize}

\begin{figure}[htbp]
  \centering
  % Segmentation combined with paging (x86): the selector indexes the segment
% descriptor table; the offset is checked against the limit and added to
% the base, giving a linear address that the page tables translate.
% Usage: % Segmentation combined with paging (x86): the selector indexes the segment
% descriptor table; the offset is checked against the limit and added to
% the base, giving a linear address that the page tables translate.
% Usage: % Segmentation combined with paging (x86): the selector indexes the segment
% descriptor table; the offset is checked against the limit and added to
% the base, giving a linear address that the page tables translate.
% Usage: \input{figures/l08-segmentation}   (width ~118 mm)
\begin{tikzpicture}[x=1mm, y=1mm, font=\scriptsize\sffamily,
  >={Stealth[length=1.5mm]},
  cell/.style={draw, minimum height=6mm, inner sep=0pt, anchor=south west},
  box/.style={draw, fill=white, align=center, inner sep=1.5pt, minimum height=9mm},
  lab/.style={font=\scriptsize\sffamily, align=center},
  hd/.style={font=\scriptsize\sffamily\bfseries, anchor=south west}]
  % logical address
  \node[hd] at (0,44.5) {\iflangja{論理アドレス}{logical address}};
  \node[cell, minimum width=16mm, fill=ln-source!16] at (0,38) {\iflangja{セレクタ}{selector}};
  \node[cell, minimum width=16mm, fill=ln-box] at (16,38) {\iflangja{オフセット}{offset}};
  % descriptor table
  \draw[fill=white] (0,28) rectangle (28,33);
  \draw[fill=ln-source!16] (0,22) rectangle (28,28) node[midway] {\iflangja{リミット}{limit}};
  \draw[fill=ln-source!16] (0,16) rectangle (28,22) node[midway] {\iflangja{ベース}{base}};
  \draw[fill=white] (0,11) rectangle (28,16);
  \draw[fill=white] (0,6) rectangle (28,11);
  \node[lab, anchor=north] at (14,5.5) {\iflangja{セグメント記述子表\\（GDT/LDT）}{segment descriptor\\table (GDT/LDT)}};
  \draw[->] (8,38) -- (8,33) node[pos=0.5, right, font=\scriptsize\sffamily, text=ln-muted] {\iflangja{索引}{index}};
  % limit check
  \node[box, minimum width=20mm] (chk) at (48,31) {\iflangja{オフセット\\$<$ リミット？}{offset\\$<$ limit?}};
  \draw[->] (32,41) -| (chk.north);
  \draw[->] (28,25) -- (34,25) |- (chk.west);
  \node[lab, anchor=west, text=ln-caution] (flt) at (64,31) {\iflangja{いいえ：フォールト}{no: fault}};
  \draw[->, ln-caution] (chk.east) -- (flt.west);
  % adder
  \node[draw, circle, inner sep=0.5pt, fill=white, font=\footnotesize] (add) at (48,13) {$+$};
  \draw[->] (chk.south) -- (add.north) node[pos=0.5, right, font=\scriptsize\sffamily] {\iflangja{はい}{yes}};
  \draw[->] (28,19) -| (40,13) -- (add.west);
  % linear -> paging -> physical
  \node[box, minimum width=18mm, fill=ln-accent!10] (lin) at (68,13) {\iflangja{リニア\\アドレス}{linear\\address}};
  \node[box, minimum width=18mm, fill=ln-accent!22] (pg) at (91,13) {\iflangja{ページ\\テーブル}{page\\tables}};
  \node[box, minimum width=16mm] (pa) at (111,13) {\iflangja{物理\\アドレス}{physical\\address}};
  \draw[->] (add.east) -- (lin.west);
  \draw[->] (lin.east) -- (pg.west);
  \draw[->] (pg.east) -- (pa.west);
\end{tikzpicture}
   (width ~118 mm)
\begin{tikzpicture}[x=1mm, y=1mm, font=\scriptsize\sffamily,
  >={Stealth[length=1.5mm]},
  cell/.style={draw, minimum height=6mm, inner sep=0pt, anchor=south west},
  box/.style={draw, fill=white, align=center, inner sep=1.5pt, minimum height=9mm},
  lab/.style={font=\scriptsize\sffamily, align=center},
  hd/.style={font=\scriptsize\sffamily\bfseries, anchor=south west}]
  % logical address
  \node[hd] at (0,44.5) {\iflangja{論理アドレス}{logical address}};
  \node[cell, minimum width=16mm, fill=ln-source!16] at (0,38) {\iflangja{セレクタ}{selector}};
  \node[cell, minimum width=16mm, fill=ln-box] at (16,38) {\iflangja{オフセット}{offset}};
  % descriptor table
  \draw[fill=white] (0,28) rectangle (28,33);
  \draw[fill=ln-source!16] (0,22) rectangle (28,28) node[midway] {\iflangja{リミット}{limit}};
  \draw[fill=ln-source!16] (0,16) rectangle (28,22) node[midway] {\iflangja{ベース}{base}};
  \draw[fill=white] (0,11) rectangle (28,16);
  \draw[fill=white] (0,6) rectangle (28,11);
  \node[lab, anchor=north] at (14,5.5) {\iflangja{セグメント記述子表\\（GDT/LDT）}{segment descriptor\\table (GDT/LDT)}};
  \draw[->] (8,38) -- (8,33) node[pos=0.5, right, font=\scriptsize\sffamily, text=ln-muted] {\iflangja{索引}{index}};
  % limit check
  \node[box, minimum width=20mm] (chk) at (48,31) {\iflangja{オフセット\\$<$ リミット？}{offset\\$<$ limit?}};
  \draw[->] (32,41) -| (chk.north);
  \draw[->] (28,25) -- (34,25) |- (chk.west);
  \node[lab, anchor=west, text=ln-caution] (flt) at (64,31) {\iflangja{いいえ：フォールト}{no: fault}};
  \draw[->, ln-caution] (chk.east) -- (flt.west);
  % adder
  \node[draw, circle, inner sep=0.5pt, fill=white, font=\footnotesize] (add) at (48,13) {$+$};
  \draw[->] (chk.south) -- (add.north) node[pos=0.5, right, font=\scriptsize\sffamily] {\iflangja{はい}{yes}};
  \draw[->] (28,19) -| (40,13) -- (add.west);
  % linear -> paging -> physical
  \node[box, minimum width=18mm, fill=ln-accent!10] (lin) at (68,13) {\iflangja{リニア\\アドレス}{linear\\address}};
  \node[box, minimum width=18mm, fill=ln-accent!22] (pg) at (91,13) {\iflangja{ページ\\テーブル}{page\\tables}};
  \node[box, minimum width=16mm] (pa) at (111,13) {\iflangja{物理\\アドレス}{physical\\address}};
  \draw[->] (add.east) -- (lin.west);
  \draw[->] (lin.east) -- (pg.west);
  \draw[->] (pg.east) -- (pa.west);
\end{tikzpicture}
   (width ~118 mm)
\begin{tikzpicture}[x=1mm, y=1mm, font=\scriptsize\sffamily,
  >={Stealth[length=1.5mm]},
  cell/.style={draw, minimum height=6mm, inner sep=0pt, anchor=south west},
  box/.style={draw, fill=white, align=center, inner sep=1.5pt, minimum height=9mm},
  lab/.style={font=\scriptsize\sffamily, align=center},
  hd/.style={font=\scriptsize\sffamily\bfseries, anchor=south west}]
  % logical address
  \node[hd] at (0,44.5) {\iflangja{論理アドレス}{logical address}};
  \node[cell, minimum width=16mm, fill=ln-source!16] at (0,38) {\iflangja{セレクタ}{selector}};
  \node[cell, minimum width=16mm, fill=ln-box] at (16,38) {\iflangja{オフセット}{offset}};
  % descriptor table
  \draw[fill=white] (0,28) rectangle (28,33);
  \draw[fill=ln-source!16] (0,22) rectangle (28,28) node[midway] {\iflangja{リミット}{limit}};
  \draw[fill=ln-source!16] (0,16) rectangle (28,22) node[midway] {\iflangja{ベース}{base}};
  \draw[fill=white] (0,11) rectangle (28,16);
  \draw[fill=white] (0,6) rectangle (28,11);
  \node[lab, anchor=north] at (14,5.5) {\iflangja{セグメント記述子表\\（GDT/LDT）}{segment descriptor\\table (GDT/LDT)}};
  \draw[->] (8,38) -- (8,33) node[pos=0.5, right, font=\scriptsize\sffamily, text=ln-muted] {\iflangja{索引}{index}};
  % limit check
  \node[box, minimum width=20mm] (chk) at (48,31) {\iflangja{オフセット\\$<$ リミット？}{offset\\$<$ limit?}};
  \draw[->] (32,41) -| (chk.north);
  \draw[->] (28,25) -- (34,25) |- (chk.west);
  \node[lab, anchor=west, text=ln-caution] (flt) at (64,31) {\iflangja{いいえ：フォールト}{no: fault}};
  \draw[->, ln-caution] (chk.east) -- (flt.west);
  % adder
  \node[draw, circle, inner sep=0.5pt, fill=white, font=\footnotesize] (add) at (48,13) {$+$};
  \draw[->] (chk.south) -- (add.north) node[pos=0.5, right, font=\scriptsize\sffamily] {\iflangja{はい}{yes}};
  \draw[->] (28,19) -| (40,13) -- (add.west);
  % linear -> paging -> physical
  \node[box, minimum width=18mm, fill=ln-accent!10] (lin) at (68,13) {\iflangja{リニア\\アドレス}{linear\\address}};
  \node[box, minimum width=18mm, fill=ln-accent!22] (pg) at (91,13) {\iflangja{ページ\\テーブル}{page\\tables}};
  \node[box, minimum width=16mm] (pa) at (111,13) {\iflangja{物理\\アドレス}{physical\\address}};
  \draw[->] (add.east) -- (lin.west);
  \draw[->] (lin.east) -- (pg.west);
  \draw[->] (pg.east) -- (pa.west);
\end{tikzpicture}

  \caption{x86 における，ページングと組み合わされたセグメンテーション．
    セレクタがセグメント記述子表の索引となる．リミット内のオフセットがベース
    に加えられ，得られたリニアアドレスがページテーブルによって変換される．}
  \label{fig:l08-segmentation}
\end{figure}

\section{ページの大きさ}
\label{sec:l08-page-size}

ページの大きさは，そのオフセットのビット数で決まる．\source{P3L2，ページの
大きさ；OSTEP ch.~20，23；Silberschatz ら ch.~9．}10 ビットのオフセットは
\SI{1}{\kibi\byte} のページを，12 ビットのオフセットは \SI{4}{\kibi\byte} の
ページを与え，後者が x86 での既定である．多くのアーキテクチャは，複数の
ページサイズを同時にサポートする．

x86 では，既定のほかに2つの大きなページサイズが使える．
\term[らーじぺーじ]{ラージページ}[large page]は \SI{2}{\mebi\byte} の大きさを
持ち，そのオフセットは 21 ビットである．オフセットが 30 ビットの
\SI{1}{\gibi\byte} のページは\term[ひゅーじぺーじ]{ヒュージページ}[huge page]と
呼ばれる．\cref{tab:l08-page-sizes}は3つの大きさを比較したものである．
（Linux は大きい方の2つのサイズのページをどちらもヒュージページと呼ぶ．）

\begin{table}[htbp]
  \centering
  \caption{x86 のページサイズと，\SI{1}{\gibi\byte} のメモリを対応付けるのに
    必要なページテーブルエントリ数．}
  \label{tab:l08-page-sizes}
  \small
  \begin{tabular}{@{}lrrr@{}}
    \toprule
    ページの大きさ & オフセット幅 & \SI{1}{\gibi\byte} のエントリ数 & 削減率 \\
    \midrule
    \SI{4}{\kibi\byte}（既定）   & 12 & 262\,144 & 1 \\
    \SI{2}{\mebi\byte}（ラージ） & 21 & 512      & $2^{9} = 512$ \\
    \SI{1}{\gibi\byte}（ヒュージ） & 30 & 1      & $2^{18} = 262\,144$ \\
    \bottomrule
  \end{tabular}
\end{table}

大きなページには2つの利点がある．必要なページテーブルエントリが少ないので
ページテーブルが小さくなり，TLB のエントリ1つがより多くのメモリを覆うので
TLB ヒットとなるアクセスが増える．欠点はメモリの無駄である．

\begin{definition}[内部断片化]
  割り当てが単位全体---ここではページ全体---に切り上げられるために，
  割り当てられていながら使われないメモリは，
  \term[ないぶだんぺんか]{内部断片化}[internal fragmentation]によって失われる．
\end{definition}

大きなページは，プロセスがその一部しか使わなくても全体が対応付けられるので，
ページが大きいほど内部断片化が無駄にしうるメモリも大きくなる．

\begin{example}
  あるプロセスが \SI{1}{\gibi\byte} のメモリを対応付ける．\SI{4}{\kibi\byte}
  のページでは 262\,144 個のエントリを要し，64 エントリの TLB はそのうち
  \SI{256}{\kibi\byte} にしか届かない．\SI{2}{\mebi\byte} のページでは 512 個
  のエントリで済み，同じ TLB が \SI{128}{\mebi\byte} に届く．しかし
  \SI{5222}{\kibi\byte} の別の領域は，\SI{4}{\kibi\byte} のページなら 1306
  ページを占めて \SI{2}{\kibi\byte} を無駄にするが，\SI{2}{\mebi\byte} の
  ページなら3ページを占めて $6144 - 5222 = \SI{922}{\kibi\byte}$ を無駄に
  する．
\end{example}

\section{メモリの割り当て}
\label{sec:l08-allocation}

アドレス変換は，仮想アドレスから物理アドレスへの対応付けを適用する．
\source{P3L2，メモリの割り当て，メモリ割り当ての課題；OSTEP ch.~17；
Silberschatz ら ch.~9--10．}ページテーブルがその対応付けを記録し，MMU が
すべてのアクセスをそれに照らして検査するが，対応付けそのもの---どのプロセスの
どのページをどの物理メモリが裏づけ，どのメモリがまだ空いているか---を決める
のは別の構成要素である．

\begin{definition}[メモリアロケータ]
  \term[めもりあろけーた]{メモリアロケータ}[memory allocator]は，仮想メモリの
  ある領域をどの物理メモリが裏づけるか，したがって仮想アドレスから物理
  アドレスへの対応付けを決定する．ページテーブルと MMU は，この対応付けを
  記録し，アクセスの正当性とその許可を検査するだけである．
\end{definition}

アロケータは2つの水準に存在する．
\begin{description}
  \item[カーネルレベルのアロケータ] プロセス制御ブロックのようなカーネルの
    状態のために，またプロセスの静的な状態---そのコード，スタック，初期化済み
    データ---のために，メモリ領域すなわちページを割り当てる．システム内の
    空きメモリの管理も行う．
  \item[ユーザレベルのアロケータ] \texttt{malloc} や \texttt{free} のような
    呼び出しを通じて，プロセスの動的な状態であるヒープを管理する．カーネルが
    プロセスにメモリの領域を与え，ユーザレベルのアロケータは，再びカーネルを
    介することなくその領域を個々の割り当てに分ける．dlmalloc，jemalloc，
    Hoard，tcmalloc などがその例である．\margin{glibc の \texttt{malloc} は
    \texttt{brk} でヒープを伸ばすが，mmap しきい値（既定で
    \SI{128}{\kibi\byte}）を超える要求には専用の \texttt{mmap} を用いる．}
\end{description}

アロケータの主要な課題は，空きメモリの断片化である．

\begin{definition}[外部断片化]
  空きメモリが連続しない穴に分かれ，総量としては十分なメモリが空いている
  にもかかわらず連続した領域の要求を満たせなくなるとき，空きメモリは
  \term[がいぶだんぺんか]{外部断片化}[external fragmentation]に悩まされる．
\end{definition}

連続した物理メモリは，例えばラージページのために，デバイスが直接アクセスする
バッファのために，そしてカーネルのデータ構造のために必要である．アロケータは
外部断片化に2つの方法で対抗する．空きメモリが連続を保つように割り当てを配置
すること，そしてメモリが解放されたときに，解放された領域を隣接する空き領域と
併合してより大きな領域にすることである．

\begin{example}
  \label{ex:l08-fragmentation}
  12 個のページフレームからなるプールが，A（3 フレーム），B（2），C（3），
  D（2）の要求を受ける．その後 A と C が解放され，E の要求が連続した 4
  フレームを必要とする（\cref{fig:l08-fragmentation}）．A と C が離して
  配置されていれば，8 フレームが空いているものの，それは 3，3，2 フレームの
  穴に分かれており，E は失敗する．A と C が隣り合って配置されていれば，
  それらを解放すると 6 フレームの連続した穴が残り，E は成功する．
\end{example}

\begin{figure}[htbp]
  \centering
  % External fragmentation: 12 page frames, requests A(3) B(2) C(3) D(2),
% then A and C are freed and E(4) needs 4 contiguous frames.
% (a) A and C placed apart: E fails.  (b) A and C placed together: E fits.
% Usage: % External fragmentation: 12 page frames, requests A(3) B(2) C(3) D(2),
% then A and C are freed and E(4) needs 4 contiguous frames.
% (a) A and C placed apart: E fails.  (b) A and C placed together: E fits.
% Usage: % External fragmentation: 12 page frames, requests A(3) B(2) C(3) D(2),
% then A and C are freed and E(4) needs 4 contiguous frames.
% (a) A and C placed apart: E fails.  (b) A and C placed together: E fits.
% Usage: \input{figures/l08-fragmentation}   (width ~118 mm)
\begin{tikzpicture}[x=1mm, y=1mm, font=\scriptsize\sffamily,
  lab/.style={font=\scriptsize\sffamily, anchor=east},
  capt/.style={font=\footnotesize\sffamily, anchor=west},
  idx/.style={font=\scriptsize\sffamily, text=ln-muted}]
  % \frow{y}{list of 12 owners}: draw a row of 12 frames of 7.5 mm
  \def\frow#1#2{%
    \foreach \own [count=\c from 0] in {#2} {
      \if\own F\def\col{white}\def\txt{}\fi
      \if\own A\def\col{ln-accent!25}\def\txt{A}\fi
      \if\own B\def\col{ln-source!20}\def\txt{B}\fi
      \if\own C\def\col{ln-tip!22}\def\txt{C}\fi
      \if\own D\def\col{ln-caution!15}\def\txt{D}\fi
      \if\own E\def\col{ln-accent!55}\def\txt{E}\fi
      \draw[fill=\col] (28+7.5*\c,#1) rectangle ++(7.5,5)
        node[midway] {\txt};
    }}
  \foreach \c in {0,...,11} { \node[idx] at (31.75+7.5*\c,53) {\c}; }
  % ---------- (a)
  \node[capt] at (0,56) {\iflangja{(a) A と C を離して配置}{(a) A and C placed apart}};
  \node[lab] at (27,47.5) {\iflangja{A--D を割り当て後}{after A--D}};
  \frow{45}{A,A,A,B,B,C,C,C,D,D,F,F}
  \node[lab] at (27,40.5) {\iflangja{A，C を解放後}{after freeing A, C}};
  \frow{38}{F,F,F,B,B,F,F,F,D,D,F,F}
  \node[lab] at (27,33.5) {\iflangja{E（4 フレーム）}{request E (4)}};
  \node[font=\scriptsize\sffamily, text=ln-caution, anchor=west] at (28,33.5)
    {\iflangja{失敗：空き 8 フレームだが連続 4 フレームがない}{fails: 8 frames free, but no 4 contiguous}};
  % ---------- (b)
  \node[capt] at (0,24) {\iflangja{(b) A と C を隣接して配置}{(b) A and C placed together}};
  \node[lab] at (27,15.5) {\iflangja{A--D を割り当て後}{after A--D}};
  \frow{13}{A,A,A,C,C,C,B,B,D,D,F,F}
  \node[lab] at (27,8.5) {\iflangja{A，C を解放後}{after freeing A, C}};
  \frow{6}{F,F,F,F,F,F,B,B,D,D,F,F}
  \node[lab] at (27,1.5) {\iflangja{E（4 フレーム）}{request E (4)}};
  \frow{-1}{E,E,E,E,F,F,B,B,D,D,F,F}
\end{tikzpicture}
   (width ~118 mm)
\begin{tikzpicture}[x=1mm, y=1mm, font=\scriptsize\sffamily,
  lab/.style={font=\scriptsize\sffamily, anchor=east},
  capt/.style={font=\footnotesize\sffamily, anchor=west},
  idx/.style={font=\scriptsize\sffamily, text=ln-muted}]
  % \frow{y}{list of 12 owners}: draw a row of 12 frames of 7.5 mm
  \def\frow#1#2{%
    \foreach \own [count=\c from 0] in {#2} {
      \if\own F\def\col{white}\def\txt{}\fi
      \if\own A\def\col{ln-accent!25}\def\txt{A}\fi
      \if\own B\def\col{ln-source!20}\def\txt{B}\fi
      \if\own C\def\col{ln-tip!22}\def\txt{C}\fi
      \if\own D\def\col{ln-caution!15}\def\txt{D}\fi
      \if\own E\def\col{ln-accent!55}\def\txt{E}\fi
      \draw[fill=\col] (28+7.5*\c,#1) rectangle ++(7.5,5)
        node[midway] {\txt};
    }}
  \foreach \c in {0,...,11} { \node[idx] at (31.75+7.5*\c,53) {\c}; }
  % ---------- (a)
  \node[capt] at (0,56) {\iflangja{(a) A と C を離して配置}{(a) A and C placed apart}};
  \node[lab] at (27,47.5) {\iflangja{A--D を割り当て後}{after A--D}};
  \frow{45}{A,A,A,B,B,C,C,C,D,D,F,F}
  \node[lab] at (27,40.5) {\iflangja{A，C を解放後}{after freeing A, C}};
  \frow{38}{F,F,F,B,B,F,F,F,D,D,F,F}
  \node[lab] at (27,33.5) {\iflangja{E（4 フレーム）}{request E (4)}};
  \node[font=\scriptsize\sffamily, text=ln-caution, anchor=west] at (28,33.5)
    {\iflangja{失敗：空き 8 フレームだが連続 4 フレームがない}{fails: 8 frames free, but no 4 contiguous}};
  % ---------- (b)
  \node[capt] at (0,24) {\iflangja{(b) A と C を隣接して配置}{(b) A and C placed together}};
  \node[lab] at (27,15.5) {\iflangja{A--D を割り当て後}{after A--D}};
  \frow{13}{A,A,A,C,C,C,B,B,D,D,F,F}
  \node[lab] at (27,8.5) {\iflangja{A，C を解放後}{after freeing A, C}};
  \frow{6}{F,F,F,F,F,F,B,B,D,D,F,F}
  \node[lab] at (27,1.5) {\iflangja{E（4 フレーム）}{request E (4)}};
  \frow{-1}{E,E,E,E,F,F,B,B,D,D,F,F}
\end{tikzpicture}
   (width ~118 mm)
\begin{tikzpicture}[x=1mm, y=1mm, font=\scriptsize\sffamily,
  lab/.style={font=\scriptsize\sffamily, anchor=east},
  capt/.style={font=\footnotesize\sffamily, anchor=west},
  idx/.style={font=\scriptsize\sffamily, text=ln-muted}]
  % \frow{y}{list of 12 owners}: draw a row of 12 frames of 7.5 mm
  \def\frow#1#2{%
    \foreach \own [count=\c from 0] in {#2} {
      \if\own F\def\col{white}\def\txt{}\fi
      \if\own A\def\col{ln-accent!25}\def\txt{A}\fi
      \if\own B\def\col{ln-source!20}\def\txt{B}\fi
      \if\own C\def\col{ln-tip!22}\def\txt{C}\fi
      \if\own D\def\col{ln-caution!15}\def\txt{D}\fi
      \if\own E\def\col{ln-accent!55}\def\txt{E}\fi
      \draw[fill=\col] (28+7.5*\c,#1) rectangle ++(7.5,5)
        node[midway] {\txt};
    }}
  \foreach \c in {0,...,11} { \node[idx] at (31.75+7.5*\c,53) {\c}; }
  % ---------- (a)
  \node[capt] at (0,56) {\iflangja{(a) A と C を離して配置}{(a) A and C placed apart}};
  \node[lab] at (27,47.5) {\iflangja{A--D を割り当て後}{after A--D}};
  \frow{45}{A,A,A,B,B,C,C,C,D,D,F,F}
  \node[lab] at (27,40.5) {\iflangja{A，C を解放後}{after freeing A, C}};
  \frow{38}{F,F,F,B,B,F,F,F,D,D,F,F}
  \node[lab] at (27,33.5) {\iflangja{E（4 フレーム）}{request E (4)}};
  \node[font=\scriptsize\sffamily, text=ln-caution, anchor=west] at (28,33.5)
    {\iflangja{失敗：空き 8 フレームだが連続 4 フレームがない}{fails: 8 frames free, but no 4 contiguous}};
  % ---------- (b)
  \node[capt] at (0,24) {\iflangja{(b) A と C を隣接して配置}{(b) A and C placed together}};
  \node[lab] at (27,15.5) {\iflangja{A--D を割り当て後}{after A--D}};
  \frow{13}{A,A,A,C,C,C,B,B,D,D,F,F}
  \node[lab] at (27,8.5) {\iflangja{A，C を解放後}{after freeing A, C}};
  \frow{6}{F,F,F,F,F,F,B,B,D,D,F,F}
  \node[lab] at (27,1.5) {\iflangja{E（4 フレーム）}{request E (4)}};
  \frow{-1}{E,E,E,E,F,F,B,B,D,D,F,F}
\end{tikzpicture}

  \caption{12 フレームのプールにおける外部断片化．同じ要求と解放の並びが，
    アロケータが A と C をどこに配置したかに応じて，(a) 3つの小さな穴を
    残すか，(b) 1つの大きな穴を残す．}
  \label{fig:l08-fragmentation}
\end{figure}

\section{Linux カーネルのアロケータ}
\label{sec:l08-linux-allocators}

Linux カーネルは2つのアロケータを組み合わせる．ページフレームのための
バディアロケータと，その上に載るカーネルオブジェクトのためのスラブアロケータ
である．\source{P3L2，Linux カーネルのアロケータ；OSTEP ch.~17；
Silberschatz ら ch.~10；Bovet と Cesati ch.~8；Bonwick \cite{bonwick1994}．}

\subsection{バディアロケータ}

\begin{definition}[バディアロケータ]
  \term[ばでぃあろけーた]{バディアロケータ}[buddy allocator]は，空きメモリを
  2の冪の大きさのブロックで管理する．要求は2の冪に切り上げられ，それを満たす
  最小の空きブロックから割り当てられる．要求された大きさのブロックができる
  まで，より大きなブロックは半分ずつに分割され，その一対をバディと呼ぶ．
  ブロックが解放され，そのバディも空いていれば，両者は併合され，併合は木を
  上へ向かって続く．
\end{definition}

それでも断片化は起こりうるが，併合はよく働き，しかも速い．フレーム $a$ から
始まる $2^{k}$ フレームのブロックのバディは，フレーム $a \oplus 2^{k}$ から
始まる．すなわち2つのアドレスは1ビットだけ異なるので，アロケータは探索せずに
バディを見つけられる．ただし2の冪への切り上げは内部断片化を生
む．\margin{Linux は 0--10 の次数，すなわち 1 個から 1024 個のフレーム
（\SI{4}{\kibi\byte} から \SI{4}{\mebi\byte}）ごとに空きリストを持
つ．\texttt{/proc/buddyinfo} が次数ごとの空きブロック数を示す．}

\begin{example}
  \label{ex:l08-buddy}
  あるバディアロケータが 16 個のページフレームを管理する．
  \cref{tab:l08-buddy}は5つの操作をたどり，\cref{fig:l08-buddy}は3番目の後の
  ブロックを示す．3 フレームの要求は 4 に，7 フレームの要求は 8 に切り上げ
  られ，それぞれ 1 フレームを無駄にする．B が解放されると，そのバディ 6--7 は
  空いているので，両者は 4--7 に併合される．4--7 のバディは
  $0 = 4 \oplus 4$ よりブロック 0--3 であり，これも空いているので併合は 0--7
  まで続く．そのバディ 8--15 は割り当て済みなので，そこで止まる．
\end{example}

\begin{table}[htbp]
  \centering
  \caption{16 個のページフレームを持つバディアロケータのトレース．}
  \label{tab:l08-buddy}
  \small
  \begin{tabular}{@{}clll@{}}
    \toprule
    ステップ & 操作 & 分割・併合 & 空きブロック \\
    \midrule
    0 & ---                      & ---                                 & 0--15 \\
    1 & A を割り当て（3 $\to$ 4） & 0--15，0--7 を分割；A $=$ 0--3      & 4--7，8--15 \\
    2 & B を割り当て（2）        & 4--7 を分割；B $=$ 4--5             & 6--7，8--15 \\
    3 & C を割り当て（7 $\to$ 8） & C $=$ 8--15                        & 6--7 \\
    4 & A を解放                 & 4--7 は分割済み，併合なし         & 0--3，6--7 \\
    5 & B を解放                 & 4--5 を 6--7 と，次に 0--3 と併合   & 0--7 \\
    \bottomrule
  \end{tabular}
\end{table}

\begin{figure}[htbp]
  \centering
  % Buddy allocator over 16 page frames after A(3 -> 4), B(2) and C(7 -> 8)
% have been allocated.  Split blocks are dashed; the unused page of a
% rounded-up allocation is lightly shaded.
% Usage: % Buddy allocator over 16 page frames after A(3 -> 4), B(2) and C(7 -> 8)
% have been allocated.  Split blocks are dashed; the unused page of a
% rounded-up allocation is lightly shaded.
% Usage: % Buddy allocator over 16 page frames after A(3 -> 4), B(2) and C(7 -> 8)
% have been allocated.  Split blocks are dashed; the unused page of a
% rounded-up allocation is lightly shaded.
% Usage: \input{figures/l08-buddy}   (width ~118 mm)
\begin{tikzpicture}[x=1mm, y=1mm, font=\scriptsize\sffamily,
  split/.style={draw, densely dashed, fill=white},
  used/.style={draw, fill=ln-accent!25},
  waste/.style={draw, fill=ln-accent!8},
  freeb/.style={draw, fill=white},
  lab/.style={font=\scriptsize\sffamily, anchor=east, text=ln-muted},
  edge/.style={ln-rule}]
  % frame f starts at x = 22 + 6 f; levels at y = 33 (16), 22 (8), 11 (4), 0 (2)
  \node[lab] at (20,36) {\iflangja{16 フレーム}{16 frames}};
  \node[lab] at (20,25) {\iflangja{8 フレーム}{8 frames}};
  \node[lab] at (20,14) {\iflangja{4 フレーム}{4 frames}};
  \node[lab] at (20,3)  {\iflangja{2 フレーム}{2 frames}};
  % level 16
  \draw[split] (22,33) rectangle (118,39) node[midway] {0--15 \iflangja{（分割）}{(split)}};
  % level 8
  \draw[split] (22,22) rectangle (70,28) node[midway] {0--7 \iflangja{（分割）}{(split)}};
  \draw[used] (70,22) rectangle (112,28) node[midway] {C (7)};
  \draw[waste] (112,22) rectangle (118,28);
  % level 4
  \draw[used] (22,11) rectangle (40,17) node[midway] {A (3)};
  \draw[waste] (40,11) rectangle (46,17);
  \draw[split] (46,11) rectangle (70,17) node[midway] {4--7 \iflangja{（分割）}{(split)}};
  % level 2
  \draw[used] (46,0) rectangle (58,6) node[midway] {B};
  \draw[freeb] (58,0) rectangle (70,6) node[midway] {\iflangja{空き}{free}};
  % tree edges
  \draw[edge] (46,33) -- (46,28);
  \draw[edge] (94,33) -- (94,28);
  \draw[edge] (34,22) -- (34,17);
  \draw[edge] (58,22) -- (58,17);
  \draw[edge] (52,11) -- (52,6);
  \draw[edge] (64,11) -- (64,6);
  % frame axis
  \draw[ln-muted] (22,-3) -- (118,-3);
  \foreach \f in {0,4,8,12,16} {
    \draw[ln-muted] (22+6*\f,-3) -- ++(0,-1);
    \node[font=\scriptsize\sffamily, text=ln-muted, anchor=north] at (22+6*\f,-4) {\f};
  }
  \node[font=\scriptsize\sffamily, text=ln-muted, anchor=north west] at (22,-8)
    {\iflangja{フレーム番号}{frame number}};
\end{tikzpicture}
   (width ~118 mm)
\begin{tikzpicture}[x=1mm, y=1mm, font=\scriptsize\sffamily,
  split/.style={draw, densely dashed, fill=white},
  used/.style={draw, fill=ln-accent!25},
  waste/.style={draw, fill=ln-accent!8},
  freeb/.style={draw, fill=white},
  lab/.style={font=\scriptsize\sffamily, anchor=east, text=ln-muted},
  edge/.style={ln-rule}]
  % frame f starts at x = 22 + 6 f; levels at y = 33 (16), 22 (8), 11 (4), 0 (2)
  \node[lab] at (20,36) {\iflangja{16 フレーム}{16 frames}};
  \node[lab] at (20,25) {\iflangja{8 フレーム}{8 frames}};
  \node[lab] at (20,14) {\iflangja{4 フレーム}{4 frames}};
  \node[lab] at (20,3)  {\iflangja{2 フレーム}{2 frames}};
  % level 16
  \draw[split] (22,33) rectangle (118,39) node[midway] {0--15 \iflangja{（分割）}{(split)}};
  % level 8
  \draw[split] (22,22) rectangle (70,28) node[midway] {0--7 \iflangja{（分割）}{(split)}};
  \draw[used] (70,22) rectangle (112,28) node[midway] {C (7)};
  \draw[waste] (112,22) rectangle (118,28);
  % level 4
  \draw[used] (22,11) rectangle (40,17) node[midway] {A (3)};
  \draw[waste] (40,11) rectangle (46,17);
  \draw[split] (46,11) rectangle (70,17) node[midway] {4--7 \iflangja{（分割）}{(split)}};
  % level 2
  \draw[used] (46,0) rectangle (58,6) node[midway] {B};
  \draw[freeb] (58,0) rectangle (70,6) node[midway] {\iflangja{空き}{free}};
  % tree edges
  \draw[edge] (46,33) -- (46,28);
  \draw[edge] (94,33) -- (94,28);
  \draw[edge] (34,22) -- (34,17);
  \draw[edge] (58,22) -- (58,17);
  \draw[edge] (52,11) -- (52,6);
  \draw[edge] (64,11) -- (64,6);
  % frame axis
  \draw[ln-muted] (22,-3) -- (118,-3);
  \foreach \f in {0,4,8,12,16} {
    \draw[ln-muted] (22+6*\f,-3) -- ++(0,-1);
    \node[font=\scriptsize\sffamily, text=ln-muted, anchor=north] at (22+6*\f,-4) {\f};
  }
  \node[font=\scriptsize\sffamily, text=ln-muted, anchor=north west] at (22,-8)
    {\iflangja{フレーム番号}{frame number}};
\end{tikzpicture}
   (width ~118 mm)
\begin{tikzpicture}[x=1mm, y=1mm, font=\scriptsize\sffamily,
  split/.style={draw, densely dashed, fill=white},
  used/.style={draw, fill=ln-accent!25},
  waste/.style={draw, fill=ln-accent!8},
  freeb/.style={draw, fill=white},
  lab/.style={font=\scriptsize\sffamily, anchor=east, text=ln-muted},
  edge/.style={ln-rule}]
  % frame f starts at x = 22 + 6 f; levels at y = 33 (16), 22 (8), 11 (4), 0 (2)
  \node[lab] at (20,36) {\iflangja{16 フレーム}{16 frames}};
  \node[lab] at (20,25) {\iflangja{8 フレーム}{8 frames}};
  \node[lab] at (20,14) {\iflangja{4 フレーム}{4 frames}};
  \node[lab] at (20,3)  {\iflangja{2 フレーム}{2 frames}};
  % level 16
  \draw[split] (22,33) rectangle (118,39) node[midway] {0--15 \iflangja{（分割）}{(split)}};
  % level 8
  \draw[split] (22,22) rectangle (70,28) node[midway] {0--7 \iflangja{（分割）}{(split)}};
  \draw[used] (70,22) rectangle (112,28) node[midway] {C (7)};
  \draw[waste] (112,22) rectangle (118,28);
  % level 4
  \draw[used] (22,11) rectangle (40,17) node[midway] {A (3)};
  \draw[waste] (40,11) rectangle (46,17);
  \draw[split] (46,11) rectangle (70,17) node[midway] {4--7 \iflangja{（分割）}{(split)}};
  % level 2
  \draw[used] (46,0) rectangle (58,6) node[midway] {B};
  \draw[freeb] (58,0) rectangle (70,6) node[midway] {\iflangja{空き}{free}};
  % tree edges
  \draw[edge] (46,33) -- (46,28);
  \draw[edge] (94,33) -- (94,28);
  \draw[edge] (34,22) -- (34,17);
  \draw[edge] (58,22) -- (58,17);
  \draw[edge] (52,11) -- (52,6);
  \draw[edge] (64,11) -- (64,6);
  % frame axis
  \draw[ln-muted] (22,-3) -- (118,-3);
  \foreach \f in {0,4,8,12,16} {
    \draw[ln-muted] (22+6*\f,-3) -- ++(0,-1);
    \node[font=\scriptsize\sffamily, text=ln-muted, anchor=north] at (22+6*\f,-4) {\f};
  }
  \node[font=\scriptsize\sffamily, text=ln-muted, anchor=north west] at (22,-8)
    {\iflangja{フレーム番号}{frame number}};
\end{tikzpicture}

  \caption{\cref{ex:l08-buddy}のバディアロケータの，ステップ 3 の後の
    ブロック．括弧内は要求されたフレーム数である．破線のブロックは分割されて
    おり，薄く塗ったフレームは内部断片化で失われている．}
  \label{fig:l08-buddy}
\end{figure}

\subsection{スラブアロケータ}

カーネルのオブジェクトの大きさは2の冪ではないので，バディアロケータはそれらに
与えるメモリの多くを無駄にしてしまう．そこで Linux カーネルは，自分の
オブジェクトをスラブアロケータ \cite{bonwick1994} で割り当て
る．\sidenote{Bonwick が最初のスラブアロケータを書いたのは Solaris~2.4 のた
めである．Linux にはこの考えの実装が3つあり---SLAB，SLUB，
小規模システム向けの SLOB---2.6.23 以降は SLUB が既定であ
る．\texttt{/proc/slabinfo} が各キャッシュとそのオブジェクトの大きさを列挙す
る．}

\begin{definition}[スラブアロケータ]
  \term[すらぶあろけーた]{スラブアロケータ}[slab allocator]は，よく使われる
  オブジェクトの型または大きさごとにキャッシュを持つ．各キャッシュは
  \emph{スラブ}，すなわち連続して割り当てられた物理メモリの領域からなり，
  スラブはオブジェクトの大きさちょうどのスロットに分割される．割り当ては，
  そのオブジェクトの型のキャッシュから空きスロットを取り，解放はそのスロットを
  キャッシュへ返す．
\end{definition}

スロットはオブジェクトの大きさちょうどなので，スラブアロケータはバディ
アロケータの切り上げによる無駄を避けられる---使われないのはスラブ末尾の余り
だけである---し，外部断片化も問題にならない．解放されたスロットは，同じ型の
次のオブジェクトにちょうど合うからである（\cref{fig:l08-slab}）．その下に
あるスラブはバディアロケータから得る．\margin{各キャッシュは CPU ごとの空き
スロットのリストも持つ．割り当ての大半は共有された状態に触れず，
ロックを必要としない（第10章）．}

\begin{figure}[htbp]
  \centering
  % Slab allocator: one cache per object type; each cache consists of slabs
% (contiguous page frames) divided into slots of exactly the object size.
% Usage: % Slab allocator: one cache per object type; each cache consists of slabs
% (contiguous page frames) divided into slots of exactly the object size.
% Usage: % Slab allocator: one cache per object type; each cache consists of slabs
% (contiguous page frames) divided into slots of exactly the object size.
% Usage: \input{figures/l08-slab}   (width ~118 mm)
\begin{tikzpicture}[x=1mm, y=1mm, font=\scriptsize\sffamily,
  usedslot/.style={draw, fill=ln-accent!25},
  freeslot/.style={draw, fill=white},
  cachebox/.style={draw, fill=ln-source!14, align=center, minimum width=24mm,
                   minimum height=10mm, inner sep=1pt},
  lab/.style={font=\scriptsize\sffamily, text=ln-muted},
  >={Stealth[length=1.4mm]}]
  % cache 1: small objects, 8 slots per slab
  \node[cachebox] (c1) at (12,27) {\iflangja{キャッシュ 1\\（小さなオブジェクト）}{cache 1\\(small objects)}};
  \foreach \k/\s in {0/usedslot,1/usedslot,2/freeslot,3/usedslot,4/usedslot,5/usedslot,6/freeslot,7/usedslot} {
    \draw[\s] (32+5*\k,24) rectangle ++(5,6);
  }
  \foreach \k/\s in {0/usedslot,1/freeslot,2/freeslot,3/freeslot,4/freeslot,5/freeslot,6/freeslot,7/freeslot} {
    \draw[\s] (78+5*\k,24) rectangle ++(5,6);
  }
  \draw[thick] (32,24) rectangle (72,30);
  \draw[thick] (78,24) rectangle (118,30);
  \draw[->] (c1.east) -- (32,27);
  \draw[->] (72,27) -- (78,27);
  % cache 2: larger objects, 3 slots per slab
  \node[cachebox] (c2) at (12,9) {\iflangja{キャッシュ 2\\（大きなオブジェクト）}{cache 2\\(large objects)}};
  \foreach \k/\s in {0/usedslot,1/usedslot,2/usedslot} {
    \draw[\s] (32+13.33*\k,6) rectangle ++(13.33,6);
  }
  \foreach \k/\s in {0/freeslot,1/usedslot,2/freeslot} {
    \draw[\s] (78+13.33*\k,6) rectangle ++(13.33,6);
  }
  \draw[thick] (32,6) rectangle (72,12);
  \draw[thick] (78,6) rectangle (118,12);
  \draw[->] (c2.east) -- (32,9);
  \draw[->] (72,9) -- (78,9);
  % labels
  \node[lab, anchor=south] at (52,30.5) {\iflangja{スラブ（連続したページフレーム）}{slab (contiguous page frames)}};
  \node[lab, anchor=south] at (98,30.5) {\iflangja{スラブ}{slab}};
  \node[lab, anchor=north] at (52,5.3) {\iflangja{スロット＝オブジェクトの大きさ}{slot size = object size}};
  \draw[usedslot] (76,-0.5) rectangle ++(3,3);
  \node[lab, anchor=west] at (79.5,1) {\iflangja{使用中}{allocated}};
  \draw[freeslot] (100,-0.5) rectangle ++(3,3);
  \node[lab, anchor=west] at (103.5,1) {\iflangja{空き}{free}};
\end{tikzpicture}
   (width ~118 mm)
\begin{tikzpicture}[x=1mm, y=1mm, font=\scriptsize\sffamily,
  usedslot/.style={draw, fill=ln-accent!25},
  freeslot/.style={draw, fill=white},
  cachebox/.style={draw, fill=ln-source!14, align=center, minimum width=24mm,
                   minimum height=10mm, inner sep=1pt},
  lab/.style={font=\scriptsize\sffamily, text=ln-muted},
  >={Stealth[length=1.4mm]}]
  % cache 1: small objects, 8 slots per slab
  \node[cachebox] (c1) at (12,27) {\iflangja{キャッシュ 1\\（小さなオブジェクト）}{cache 1\\(small objects)}};
  \foreach \k/\s in {0/usedslot,1/usedslot,2/freeslot,3/usedslot,4/usedslot,5/usedslot,6/freeslot,7/usedslot} {
    \draw[\s] (32+5*\k,24) rectangle ++(5,6);
  }
  \foreach \k/\s in {0/usedslot,1/freeslot,2/freeslot,3/freeslot,4/freeslot,5/freeslot,6/freeslot,7/freeslot} {
    \draw[\s] (78+5*\k,24) rectangle ++(5,6);
  }
  \draw[thick] (32,24) rectangle (72,30);
  \draw[thick] (78,24) rectangle (118,30);
  \draw[->] (c1.east) -- (32,27);
  \draw[->] (72,27) -- (78,27);
  % cache 2: larger objects, 3 slots per slab
  \node[cachebox] (c2) at (12,9) {\iflangja{キャッシュ 2\\（大きなオブジェクト）}{cache 2\\(large objects)}};
  \foreach \k/\s in {0/usedslot,1/usedslot,2/usedslot} {
    \draw[\s] (32+13.33*\k,6) rectangle ++(13.33,6);
  }
  \foreach \k/\s in {0/freeslot,1/usedslot,2/freeslot} {
    \draw[\s] (78+13.33*\k,6) rectangle ++(13.33,6);
  }
  \draw[thick] (32,6) rectangle (72,12);
  \draw[thick] (78,6) rectangle (118,12);
  \draw[->] (c2.east) -- (32,9);
  \draw[->] (72,9) -- (78,9);
  % labels
  \node[lab, anchor=south] at (52,30.5) {\iflangja{スラブ（連続したページフレーム）}{slab (contiguous page frames)}};
  \node[lab, anchor=south] at (98,30.5) {\iflangja{スラブ}{slab}};
  \node[lab, anchor=north] at (52,5.3) {\iflangja{スロット＝オブジェクトの大きさ}{slot size = object size}};
  \draw[usedslot] (76,-0.5) rectangle ++(3,3);
  \node[lab, anchor=west] at (79.5,1) {\iflangja{使用中}{allocated}};
  \draw[freeslot] (100,-0.5) rectangle ++(3,3);
  \node[lab, anchor=west] at (103.5,1) {\iflangja{空き}{free}};
\end{tikzpicture}
   (width ~118 mm)
\begin{tikzpicture}[x=1mm, y=1mm, font=\scriptsize\sffamily,
  usedslot/.style={draw, fill=ln-accent!25},
  freeslot/.style={draw, fill=white},
  cachebox/.style={draw, fill=ln-source!14, align=center, minimum width=24mm,
                   minimum height=10mm, inner sep=1pt},
  lab/.style={font=\scriptsize\sffamily, text=ln-muted},
  >={Stealth[length=1.4mm]}]
  % cache 1: small objects, 8 slots per slab
  \node[cachebox] (c1) at (12,27) {\iflangja{キャッシュ 1\\（小さなオブジェクト）}{cache 1\\(small objects)}};
  \foreach \k/\s in {0/usedslot,1/usedslot,2/freeslot,3/usedslot,4/usedslot,5/usedslot,6/freeslot,7/usedslot} {
    \draw[\s] (32+5*\k,24) rectangle ++(5,6);
  }
  \foreach \k/\s in {0/usedslot,1/freeslot,2/freeslot,3/freeslot,4/freeslot,5/freeslot,6/freeslot,7/freeslot} {
    \draw[\s] (78+5*\k,24) rectangle ++(5,6);
  }
  \draw[thick] (32,24) rectangle (72,30);
  \draw[thick] (78,24) rectangle (118,30);
  \draw[->] (c1.east) -- (32,27);
  \draw[->] (72,27) -- (78,27);
  % cache 2: larger objects, 3 slots per slab
  \node[cachebox] (c2) at (12,9) {\iflangja{キャッシュ 2\\（大きなオブジェクト）}{cache 2\\(large objects)}};
  \foreach \k/\s in {0/usedslot,1/usedslot,2/usedslot} {
    \draw[\s] (32+13.33*\k,6) rectangle ++(13.33,6);
  }
  \foreach \k/\s in {0/freeslot,1/usedslot,2/freeslot} {
    \draw[\s] (78+13.33*\k,6) rectangle ++(13.33,6);
  }
  \draw[thick] (32,6) rectangle (72,12);
  \draw[thick] (78,6) rectangle (118,12);
  \draw[->] (c2.east) -- (32,9);
  \draw[->] (72,9) -- (78,9);
  % labels
  \node[lab, anchor=south] at (52,30.5) {\iflangja{スラブ（連続したページフレーム）}{slab (contiguous page frames)}};
  \node[lab, anchor=south] at (98,30.5) {\iflangja{スラブ}{slab}};
  \node[lab, anchor=north] at (52,5.3) {\iflangja{スロット＝オブジェクトの大きさ}{slot size = object size}};
  \draw[usedslot] (76,-0.5) rectangle ++(3,3);
  \node[lab, anchor=west] at (79.5,1) {\iflangja{使用中}{allocated}};
  \draw[freeslot] (100,-0.5) rectangle ++(3,3);
  \node[lab, anchor=west] at (103.5,1) {\iflangja{空き}{free}};
\end{tikzpicture}

  \caption{小さなオブジェクト用と大きなオブジェクト用のキャッシュを持つ
    スラブアロケータ．どのスラブもオブジェクトの大きさのスロットに分割される．}
  \label{fig:l08-slab}
\end{figure}

\begin{example}
  あるカーネルオブジェクトが 1200 バイトを占めるとする．バディアロケータから
  個別に割り当てると，各オブジェクトは少なくとも \SI{4}{\kibi\byte} の 1
  ページを受け取り，2896 バイトを無駄にする．1 ページのスラブはこのような
  オブジェクトを $\lfloor 4096/1200 \rfloor = 3$ 個保持し，
  $4096 - 3600 = 496$ バイトを未使用のまま残す．オブジェクト当たり約 165
  バイトである．
\end{example}

\needspace{9\baselineskip}
\begin{keypoint}
  Linux はページフレームをバディアロケータで割り当てる．その2の冪のブロック
  は探索なしに見つけて併合できるが，要求は切り上げられる．カーネル
  オブジェクトはスラブアロケータで割り当てる．そのスロットはオブジェクトの
  大きさを持ち，同じ型の次のオブジェクトに再利用される．
\end{keypoint}

\section{デマンドページング}
\label{sec:l08-demand}

仮想メモリは物理メモリよりはるかに大きいので，仮想ページが常にページフレーム
に存在しているとは限らない．\source{P3L2，デマンドページング；OSTEP
ch.~21；Silberschatz ら ch.~10．}ページフレームの内容は二次記憶へ退避され，
そこから復元される．

\begin{definition}[スワッピングとデマンドページング]
  ページの内容を物理メモリと二次記憶との間で移動することを
  \term[すわっぴんぐ]{スワッピング}[swapping]と呼ぶ．ページは
  \term[すわっぷりょういき]{スワップ領域}[swap partition]に格納される．これは
  ディスク，フラッシュメモリ，あるいは遠隔マシンのメモリの一部である．
  \term[でまんどぺーじんぐ]{デマンドページング}[demand paging]では，ページが
  メモリへスワップインされるのは，プロセスがそれにアクセスしたときだけである．
\end{definition}

プロセスがメモリにないページにアクセスすると，次の手順が踏まれる
（\cref{fig:l08-page-fault}）．
\begin{enumerate}
  \item プロセスがメモリアドレスを参照し，MMU はその PTE の有効ビットが
    クリアされていることを見つける．
  \item MMU がページフォールトを起こし，オペレーティングシステムへトラップ
    する．
  \item オペレーティングシステムは，そのアクセスが正当であること，および
    ページがスワップ領域のどこに格納されているかを判断する．
  \item I/O 操作を用いて，スワップ領域からページを空いているページフレームへ
    読み込む．\margin{後の例の \SI{8}{\milli\second} は回転ディスクの値であ
    る．NVMe SSD は数十マイクロ秒でページを読むので，許容できるフォールト率
    は2桁高くなる．}
  \item 新しい PFN で PTE を更新し，有効ビットを立てる．
  \item プログラムカウンタをフォールトした命令に戻してプロセスを続行させる．
    参照は再び実行され，今度は成功する．
\end{enumerate}
スワップインの後にページが占めるページフレームは，一般に以前占めていたものと
は異なる．空いているページフレームがなければ，オペレーティングシステムはまず
1つを解放しなければならない（\cref{sec:l08-replacement}）．

\begin{figure}[htbp]
  \centering
  % Demand paging: the steps taken when a process references a page that is
% in the swap partition.
% Usage: % Demand paging: the steps taken when a process references a page that is
% in the swap partition.
% Usage: % Demand paging: the steps taken when a process references a page that is
% in the swap partition.
% Usage: \input{figures/l08-page-fault}   (width ~118 mm)
\begin{tikzpicture}[x=1mm, y=1mm, font=\scriptsize\sffamily,
  >={Stealth[length=1.5mm]},
  box/.style={draw, fill=white, align=center, inner sep=1pt},
  stp/.style={draw, circle, fill=ln-accent, text=white, inner sep=0.4pt,
               font=\scriptsize\sffamily\bfseries, minimum size=4mm},
  lab/.style={font=\scriptsize\sffamily, align=left},
  hd/.style={font=\scriptsize\sffamily\bfseries, anchor=south}]
  % process
  \node[box, minimum width=22mm, minimum height=13mm] (proc) at (11,38.5)
    {\iflangja{プロセス}{process}\\[1pt]\texttt{x = a[i];}};
  % page table
  \node[hd, anchor=north] at (46,25.5) {\iflangja{ページテーブル}{page table}};
  \foreach \r in {0,...,4} { \draw[fill=white] (38,26+5*\r) rectangle ++(16,5); }
  \draw[fill=ln-caution!15] (38,36) rectangle ++(16,5) node[midway] {\iflangja{無効}{invalid}};
  % OS
  \node[box, minimum width=44mm, minimum height=8mm, fill=ln-source!12] (os) at (46,62)
    {\iflangja{OS（ページフォールトハンドラ）}{OS (page fault handler)}};
  % physical memory
  \node[hd] at (104,50.5) {\iflangja{物理メモリ}{physical memory}};
  \foreach \r in {0,...,4} { \draw[fill=ln-box] (92,26+5*\r) rectangle ++(24,5); }
  \draw[fill=ln-accent!25] (92,26) rectangle ++(24,5) node[midway] {\iflangja{空きフレーム}{free frame}};
  % swap partition
  \node[box, minimum width=28mm, minimum height=8mm, fill=ln-box, rounded corners=3pt] (swap) at (104,6)
    {\iflangja{スワップ領域}{swap partition}};
  % 1 reference
  \draw[->] (proc.east) -- (38,38.5);
  \node[stp] at (30,41) {1};
  % 2 trap
  \draw[->] (42,51) -- (42,58);
  \node[stp] at (39,54.5) {2};
  \node[lab, anchor=east] at (37,54.5) {\iflangja{トラップ}{trap}};
  % 3 locate page in swap
  \draw[->] (68.5,62) -- (80,62) |- (swap.west);
  \node[stp] at (77,58) {3};
  \node[lab, anchor=east] at (79,46) {\iflangja{スワップ上の\\位置を特定}{locate page\\in swap}};
  % 4 read page into free frame
  \draw[->, very thick, ln-accent] (104,10) -- (104,26);
  \node[stp] at (107.5,17) {4};
  \node[lab, anchor=east] at (101,17) {\iflangja{読み込み}{read page in}};
  % 5 update page table entry
  \draw[->] (50,58) -- (50,51);
  \node[stp] at (53,54.5) {5};
  \node[lab, anchor=west] at (55,54.5) {\iflangja{PTE を更新}{update PTE}};
  \draw[->, densely dashed] (54,38.5) -- (92,28.5);
  % 6 restart
  \draw[->] (os.west) -| (11,44.5);
  \node[stp] at (14,58) {6};
  \node[lab, anchor=west] at (16,54) {\iflangja{命令を\\再実行}{restart}};
\end{tikzpicture}
   (width ~118 mm)
\begin{tikzpicture}[x=1mm, y=1mm, font=\scriptsize\sffamily,
  >={Stealth[length=1.5mm]},
  box/.style={draw, fill=white, align=center, inner sep=1pt},
  stp/.style={draw, circle, fill=ln-accent, text=white, inner sep=0.4pt,
               font=\scriptsize\sffamily\bfseries, minimum size=4mm},
  lab/.style={font=\scriptsize\sffamily, align=left},
  hd/.style={font=\scriptsize\sffamily\bfseries, anchor=south}]
  % process
  \node[box, minimum width=22mm, minimum height=13mm] (proc) at (11,38.5)
    {\iflangja{プロセス}{process}\\[1pt]\texttt{x = a[i];}};
  % page table
  \node[hd, anchor=north] at (46,25.5) {\iflangja{ページテーブル}{page table}};
  \foreach \r in {0,...,4} { \draw[fill=white] (38,26+5*\r) rectangle ++(16,5); }
  \draw[fill=ln-caution!15] (38,36) rectangle ++(16,5) node[midway] {\iflangja{無効}{invalid}};
  % OS
  \node[box, minimum width=44mm, minimum height=8mm, fill=ln-source!12] (os) at (46,62)
    {\iflangja{OS（ページフォールトハンドラ）}{OS (page fault handler)}};
  % physical memory
  \node[hd] at (104,50.5) {\iflangja{物理メモリ}{physical memory}};
  \foreach \r in {0,...,4} { \draw[fill=ln-box] (92,26+5*\r) rectangle ++(24,5); }
  \draw[fill=ln-accent!25] (92,26) rectangle ++(24,5) node[midway] {\iflangja{空きフレーム}{free frame}};
  % swap partition
  \node[box, minimum width=28mm, minimum height=8mm, fill=ln-box, rounded corners=3pt] (swap) at (104,6)
    {\iflangja{スワップ領域}{swap partition}};
  % 1 reference
  \draw[->] (proc.east) -- (38,38.5);
  \node[stp] at (30,41) {1};
  % 2 trap
  \draw[->] (42,51) -- (42,58);
  \node[stp] at (39,54.5) {2};
  \node[lab, anchor=east] at (37,54.5) {\iflangja{トラップ}{trap}};
  % 3 locate page in swap
  \draw[->] (68.5,62) -- (80,62) |- (swap.west);
  \node[stp] at (77,58) {3};
  \node[lab, anchor=east] at (79,46) {\iflangja{スワップ上の\\位置を特定}{locate page\\in swap}};
  % 4 read page into free frame
  \draw[->, very thick, ln-accent] (104,10) -- (104,26);
  \node[stp] at (107.5,17) {4};
  \node[lab, anchor=east] at (101,17) {\iflangja{読み込み}{read page in}};
  % 5 update page table entry
  \draw[->] (50,58) -- (50,51);
  \node[stp] at (53,54.5) {5};
  \node[lab, anchor=west] at (55,54.5) {\iflangja{PTE を更新}{update PTE}};
  \draw[->, densely dashed] (54,38.5) -- (92,28.5);
  % 6 restart
  \draw[->] (os.west) -| (11,44.5);
  \node[stp] at (14,58) {6};
  \node[lab, anchor=west] at (16,54) {\iflangja{命令を\\再実行}{restart}};
\end{tikzpicture}
   (width ~118 mm)
\begin{tikzpicture}[x=1mm, y=1mm, font=\scriptsize\sffamily,
  >={Stealth[length=1.5mm]},
  box/.style={draw, fill=white, align=center, inner sep=1pt},
  stp/.style={draw, circle, fill=ln-accent, text=white, inner sep=0.4pt,
               font=\scriptsize\sffamily\bfseries, minimum size=4mm},
  lab/.style={font=\scriptsize\sffamily, align=left},
  hd/.style={font=\scriptsize\sffamily\bfseries, anchor=south}]
  % process
  \node[box, minimum width=22mm, minimum height=13mm] (proc) at (11,38.5)
    {\iflangja{プロセス}{process}\\[1pt]\texttt{x = a[i];}};
  % page table
  \node[hd, anchor=north] at (46,25.5) {\iflangja{ページテーブル}{page table}};
  \foreach \r in {0,...,4} { \draw[fill=white] (38,26+5*\r) rectangle ++(16,5); }
  \draw[fill=ln-caution!15] (38,36) rectangle ++(16,5) node[midway] {\iflangja{無効}{invalid}};
  % OS
  \node[box, minimum width=44mm, minimum height=8mm, fill=ln-source!12] (os) at (46,62)
    {\iflangja{OS（ページフォールトハンドラ）}{OS (page fault handler)}};
  % physical memory
  \node[hd] at (104,50.5) {\iflangja{物理メモリ}{physical memory}};
  \foreach \r in {0,...,4} { \draw[fill=ln-box] (92,26+5*\r) rectangle ++(24,5); }
  \draw[fill=ln-accent!25] (92,26) rectangle ++(24,5) node[midway] {\iflangja{空きフレーム}{free frame}};
  % swap partition
  \node[box, minimum width=28mm, minimum height=8mm, fill=ln-box, rounded corners=3pt] (swap) at (104,6)
    {\iflangja{スワップ領域}{swap partition}};
  % 1 reference
  \draw[->] (proc.east) -- (38,38.5);
  \node[stp] at (30,41) {1};
  % 2 trap
  \draw[->] (42,51) -- (42,58);
  \node[stp] at (39,54.5) {2};
  \node[lab, anchor=east] at (37,54.5) {\iflangja{トラップ}{trap}};
  % 3 locate page in swap
  \draw[->] (68.5,62) -- (80,62) |- (swap.west);
  \node[stp] at (77,58) {3};
  \node[lab, anchor=east] at (79,46) {\iflangja{スワップ上の\\位置を特定}{locate page\\in swap}};
  % 4 read page into free frame
  \draw[->, very thick, ln-accent] (104,10) -- (104,26);
  \node[stp] at (107.5,17) {4};
  \node[lab, anchor=east] at (101,17) {\iflangja{読み込み}{read page in}};
  % 5 update page table entry
  \draw[->] (50,58) -- (50,51);
  \node[stp] at (53,54.5) {5};
  \node[lab, anchor=west] at (55,54.5) {\iflangja{PTE を更新}{update PTE}};
  \draw[->, densely dashed] (54,38.5) -- (92,28.5);
  % 6 restart
  \draw[->] (os.west) -| (11,44.5);
  \node[stp] at (14,58) {6};
  \node[lab, anchor=west] at (16,54) {\iflangja{命令を\\再実行}{restart}};
\end{tikzpicture}

  \caption{デマンドページング．プロセスがスワップ領域にあるページを参照した
    ときに踏まれる手順．}
  \label{fig:l08-page-fault}
\end{figure}

スワップしてはならないページもある．\term[ぴんどめぺーじ]{ピン留めページ}[pinned page]
とは，スワッピングが禁止され，物理メモリに留まり続けるページである．例えば，
ダイレクトメモリアクセス（第11章）の際にデバイスが物理アドレスで直接
アクセスするページはピン留めされる．

\begin{example}
  メモリアクセスに \SI{100}{\nano\second}，ディスクからページを読むページ
  フォールトの処理に \SI{8}{\milli\second}，すなわち $8 \cdot 10^{6}$ ナノ秒，
  つまり 80\,000 倍の時間がかかるとする．アクセスのうち割合 $p$ がフォールト
  するなら，平均のアクセスは $(1-p) \cdot 100 + p \cdot 8 \cdot 10^{6}$ ナノ秒
  かかる．$p = 10^{-4}$ ではこれは約 \SI{900}{\nano\second} であり，9 倍遅い．
  \SI{100}{\nano\second} の \SI{10}{\percent} 以内に収めるには，$p$ は
  $1.25 \cdot 10^{-6}$，すなわち 800\,000 回のアクセスに 1 回のフォールトを
  超えてはならない．
\end{example}

\section{ページ置換}
\label{sec:l08-replacement}

物理メモリが不足すると，オペレーティングシステムはページをスワップアウトして
ページフレームを解放しなければならない．\source{P3L2，物理メモリの解放；
OSTEP ch.~22；Silberschatz ら ch.~10；Bovet と Cesati ch.~17．}スワップ
アウトされたページがすぐにまた必要になると，ページフォールトと I/O 操作の
代価がかかる．\cref{sec:l08-demand}の例が示したように，これは高くつく．
したがってオペレーティングシステムは，スワップアウトするページを注意深く
選ばなければならない．この選択を\term[ぺーじちかん]{ページ置換}[page replacement]と
呼び，そこには2つの問いがある．いつページをスワップアウトすべきか，そして
どのページをスワップアウトすべきかである．

ページをスワップアウトするのは
\term[ぺーじあうとでーもん]{ページアウトデーモン}[page-out daemon]である．
これは，メモリの使用量がしきい値である\emph{高水位}を超えたときに走るカーネル
スレッドである．アプリケーションをできるだけ妨げないよう，CPU の使用率が
しきい値である\emph{低水位}を下回っている間に走ることが望ましい．

理想的な犠牲ページは，将来使われないページである．将来は分からないので，オペレー
ティングシステムはアクセスの履歴からそれを予測する．
\begin{itemize}
  \item \term{LRU}[least recently used]%
    \index{least recently used|see{LRU}}（最長時間未使用）方針は，最後の
    アクセスが最も過去にあるページをスワップアウトする．最近使われていない
    ページは近いうちにも使われないだろう，という仮定に基づく．PTE の
    アクセスビットが，ビットが最後にクリアされてからそのページが参照されたか
    どうかを教える．
  \item 書き出す必要のないページは置換の代価が小さい．ダーティビットは，
    ページが変更されたかどうかを教える．すでにディスクに格納されている
    クリーンなページの内容は，単に捨ててよい．
  \item ピン留めページやカーネルの特定のページのように，メモリに留まらなければ
    ならないページは決して選ばれない．
\end{itemize}
Linux は，回収が始まり，また止まるタイミングを調整するパラメータを
備え，\caution{Linux では2つの水位はいずれも空きメモリに対するしきい値で
あり，低い方を下回ると回収が始まり，高い方を上回ると止まる．}ページを回収
可能なページやスワップ可能なページといった種類に分類し，ページに二度目の
機会を与える LRU の変種を用いる．

\begin{definition}[セカンドチャンスアルゴリズム]
  \term[せかんどちゃんすあるごりずむ]{セカンドチャンスアルゴリズム}[second chance algorithm]は，
  アクセスビットを用いて LRU を近似する．ページフレームを巡回順に走査する．
  アクセスビットが立っているページには二度目の機会が与えられる．そのビットは
  クリアされ，走査は次へ進む．アクセスビットがクリアされているページ，すなわち
  前回の走査がそのビットをクリアしてから参照されていないページが，スワップ
  アウトされる．
\end{definition}

\begin{example}
  \label{ex:l08-second-chance}
  5 個のページフレームがページ A--E を保持し，そのアクセスビットは 1，0，1，
  1，0 で，走査は A から始まる．\cref{tab:l08-second-chance}は2回の置換を
  たどる．2回目の走査は F を受け取ったフレームの次から再開するので，C，D，E
  だけを調べる．C と D には二度目の機会が与えられ，ビットがクリアされている E
  がスワップアウトされる．A はまったく調べられない．最初の走査がそのビットを
  クリアした後にプロセスが再び A を参照したので，次に走査が A に達したとき，
  A には二度目の機会が与えられる．そして，ステップ 3 でビットがクリアされ，
  その後参照されていない C が犠牲となる．
\end{example}

\begin{table}[htbp]
  \centering
  \caption{セカンドチャンスアルゴリズムのトレース．括弧内は読み込まれる
    ページである．アクセスビットは各ステップの後の値を走査順に並べてある．
    新たに読み込まれたページはビットが立っている．}
  \label{tab:l08-second-chance}
  \small
  \begin{tabular}{@{}clll@{}}
    \toprule
    ステップ & 事象 & 調べたページ & アクセスビット \\
    \midrule
    0 & ---            & ---                           & A1 B0 C1 D1 E0 \\
    1 & フォールト（F） & A クリア，B スワップアウト     & A0 F1 C1 D1 E0 \\
    2 & A と D を参照   & ---                           & A1 F1 C1 D1 E0 \\
    3 & フォールト（G） & C，D クリア，E スワップアウト  & A1 F1 C0 D0 G1 \\
    \bottomrule
  \end{tabular}
\end{table}

\needspace{10\baselineskip}
\begin{keypoint}
  デマンドページングは，最初にそのページを必要とするページフォールトの際に
  ページを持ち込み，ページアウトデーモンはメモリが尽きる前にフレームを解放
  する．どのページが去るかは，アクセスビットを読んでクリアするセカンド
  チャンスアルゴリズムのような LRU の近似で決められる．ダーティビットは，
  犠牲となるページをそもそも書き出す必要があるかどうかを教える．
\end{keypoint}

\section{コピーオンライト}
\label{sec:l08-cow}

MMU はアドレスを変換し，アクセスを追跡し，保護を強制する．これらの能力は，
他のサービスや最適化を作るのにも役立つ．\source{P3L2，コピーオンライト；
OSTEP ch.~23；Silberschatz ら ch.~10．}その1つは，プロセスが生成されるときの
不要なコピーを避ける．

\texttt{fork}\index{fork@\texttt{fork}} で生成されたプロセスは，親のアドレス
空間全体のコピーを受け取る（第2章）．\margin{ページテーブル自体はいずれにせ
よコピーされる．\SI{4}{\kibi\byte} のページで \SI{1}{\gibi\byte} なら
8 バイトのエントリ $2^{18}$ 個，\SI{2}{\mebi\byte} である．}そのページの多くは，
コードのように静的
で決して変わらず，しかも子はしばしば生成の直後に \texttt{exec} を呼ぶので，
コピーの大半は無駄になってしまう．

\begin{definition}[コピーオンライト]
  \term[こぴーおんらいと]{コピーオンライト}[copy-on-write]%
  \index{COW|see{コピーオンライト}}（COW）では，オペレーティングシステムは
  プロセスを生成するときにアドレス空間をコピーしない．新しいプロセスの仮想
  ページを元のプロセスのページフレームに対応付け，両方のプロセスでそれらの
  フレームを書き込み禁止にする．そのようなページへの書き込みはページ
  フォールトを起こし，そのときになって初めてオペレーティングシステムはその
  ページをコピーする．
\end{definition}

両方のプロセスが共有ページを読むだけである限り，メモリも，コピーの時間も
費やされない．どちらかがページに書き込むと，保護フォールトによってオペレー
ティングシステムはそれが COW ページであることを認識する．そしてこの1ページ
だけを新しいページフレームへコピーし，書き込んだプロセスのページを書き込み
許可付きでそのコピーに対応付け，書き込みをやり直させる（\cref{fig:l08-cow}）．
コピーの代価は，実際に変更されたページに対してだけ支払われる．\caution{節約
は無条件ではない．最初の書き込みはいずれもカーネルへのトラップを伴い，
その代価はマイクロ秒の程度である．ほとんどすべてのページに書き込む子では，
先にアドレス空間をコピーしてしまう方が速い．}

\begin{figure}[htbp]
  \centering
  % Copy-on-write after fork: (a) parent and child map the same write-
% protected frames; (b) the child's write to page 1 faults, and the OS
% copies only that page.
% Usage: % Copy-on-write after fork: (a) parent and child map the same write-
% protected frames; (b) the child's write to page 1 faults, and the OS
% copies only that page.
% Usage: % Copy-on-write after fork: (a) parent and child map the same write-
% protected frames; (b) the child's write to page 1 faults, and the OS
% copies only that page.
% Usage: \input{figures/l08-cow}   (width ~118 mm)
\begin{tikzpicture}[x=1mm, y=1mm, font=\scriptsize\sffamily,
  >={Stealth[length=1.4mm]},
  ro/.style={draw, fill=ln-caution!12},
  rw/.style={draw, fill=ln-accent!22},
  oth/.style={draw, fill=ln-box},
  hd/.style={font=\scriptsize\sffamily\bfseries, anchor=base},
  capt/.style={font=\footnotesize\sffamily},
  lab/.style={font=\scriptsize\sffamily}]
  % \panel{x0}{child page-1 frame}{copy?}
  \def\panel#1#2{%
    \begin{scope}[shift={(#1,0)}]
      \node[hd] at (6,37.5) {\iflangja{親}{parent}};
      \node[hd] at (28,37.5) {\iflangja{メモリ}{memory}};
      \node[hd] at (50,37.5) {\iflangja{子}{child}};
      % frames 0..4, height 7, centres 3.5 + 7 i
      \foreach \i in {0,...,4} { \draw[oth] (22,7*\i) rectangle ++(12,7); }
      \foreach \i/\n in {3/0, 1/1, 4/2} {
        \draw[ro] (22,7*\i) rectangle ++(12,7) node[midway, lab] {\iflangja{ページ}{page} \n};
      }
      % parent pages 0..2 (centres 10, 18, 26)
      \foreach \j/\i in {0/3, 1/1, 2/4} {
        \draw[ro] (0,6+8*\j) rectangle ++(12,8) node[midway, lab] {\j};
        \draw[->] (12,10+8*\j) -- (22,3.5+7*\i);
      }
      #2
    \end{scope}}
  % (a) after fork: child maps the same frames
  \panel{0}{%
    \foreach \j/\i in {0/3, 1/1, 2/4} {
      \draw[ro] (44,6+8*\j) rectangle ++(12,8) node[midway, lab] {\j};
      \draw[->] (44,10+8*\j) -- (34,3.5+7*\i);
    }}
  \node[capt] at (28,-5) {\iflangja{(a) \texttt{fork} 直後}{(a) after \texttt{fork}}};
  % (b) after the child writes page 1
  \panel{62}{%
    \foreach \j/\i in {0/3, 2/4} {
      \draw[ro] (44,6+8*\j) rectangle ++(12,8) node[midway, lab] {\j};
      \draw[->] (44,10+8*\j) -- (34,3.5+7*\i);
    }
    \draw[rw] (44,14) rectangle ++(12,8) node[midway, lab] {1};
    \draw[rw] (22,0) rectangle ++(12,7) node[midway, lab] {\iflangja{コピー}{copy}};
    \draw[->, very thick, ln-accent] (44,18) -- (34,3.5);}
  \node[capt] at (90,-5) {\iflangja{(b) 子がページ 1 に書き込んだ後}{(b) after the child writes page 1}};
  % legend
  \draw[ro] (0,-12) rectangle ++(4,3);
  \node[lab, anchor=west] at (5,-10.5) {\iflangja{書き込み禁止}{write-protected}};
  \draw[rw] (32,-12) rectangle ++(4,3);
  \node[lab, anchor=west] at (37,-10.5) {\iflangja{書き込み可能}{writable}};
\end{tikzpicture}
   (width ~118 mm)
\begin{tikzpicture}[x=1mm, y=1mm, font=\scriptsize\sffamily,
  >={Stealth[length=1.4mm]},
  ro/.style={draw, fill=ln-caution!12},
  rw/.style={draw, fill=ln-accent!22},
  oth/.style={draw, fill=ln-box},
  hd/.style={font=\scriptsize\sffamily\bfseries, anchor=base},
  capt/.style={font=\footnotesize\sffamily},
  lab/.style={font=\scriptsize\sffamily}]
  % \panel{x0}{child page-1 frame}{copy?}
  \def\panel#1#2{%
    \begin{scope}[shift={(#1,0)}]
      \node[hd] at (6,37.5) {\iflangja{親}{parent}};
      \node[hd] at (28,37.5) {\iflangja{メモリ}{memory}};
      \node[hd] at (50,37.5) {\iflangja{子}{child}};
      % frames 0..4, height 7, centres 3.5 + 7 i
      \foreach \i in {0,...,4} { \draw[oth] (22,7*\i) rectangle ++(12,7); }
      \foreach \i/\n in {3/0, 1/1, 4/2} {
        \draw[ro] (22,7*\i) rectangle ++(12,7) node[midway, lab] {\iflangja{ページ}{page} \n};
      }
      % parent pages 0..2 (centres 10, 18, 26)
      \foreach \j/\i in {0/3, 1/1, 2/4} {
        \draw[ro] (0,6+8*\j) rectangle ++(12,8) node[midway, lab] {\j};
        \draw[->] (12,10+8*\j) -- (22,3.5+7*\i);
      }
      #2
    \end{scope}}
  % (a) after fork: child maps the same frames
  \panel{0}{%
    \foreach \j/\i in {0/3, 1/1, 2/4} {
      \draw[ro] (44,6+8*\j) rectangle ++(12,8) node[midway, lab] {\j};
      \draw[->] (44,10+8*\j) -- (34,3.5+7*\i);
    }}
  \node[capt] at (28,-5) {\iflangja{(a) \texttt{fork} 直後}{(a) after \texttt{fork}}};
  % (b) after the child writes page 1
  \panel{62}{%
    \foreach \j/\i in {0/3, 2/4} {
      \draw[ro] (44,6+8*\j) rectangle ++(12,8) node[midway, lab] {\j};
      \draw[->] (44,10+8*\j) -- (34,3.5+7*\i);
    }
    \draw[rw] (44,14) rectangle ++(12,8) node[midway, lab] {1};
    \draw[rw] (22,0) rectangle ++(12,7) node[midway, lab] {\iflangja{コピー}{copy}};
    \draw[->, very thick, ln-accent] (44,18) -- (34,3.5);}
  \node[capt] at (90,-5) {\iflangja{(b) 子がページ 1 に書き込んだ後}{(b) after the child writes page 1}};
  % legend
  \draw[ro] (0,-12) rectangle ++(4,3);
  \node[lab, anchor=west] at (5,-10.5) {\iflangja{書き込み禁止}{write-protected}};
  \draw[rw] (32,-12) rectangle ++(4,3);
  \node[lab, anchor=west] at (37,-10.5) {\iflangja{書き込み可能}{writable}};
\end{tikzpicture}
   (width ~118 mm)
\begin{tikzpicture}[x=1mm, y=1mm, font=\scriptsize\sffamily,
  >={Stealth[length=1.4mm]},
  ro/.style={draw, fill=ln-caution!12},
  rw/.style={draw, fill=ln-accent!22},
  oth/.style={draw, fill=ln-box},
  hd/.style={font=\scriptsize\sffamily\bfseries, anchor=base},
  capt/.style={font=\footnotesize\sffamily},
  lab/.style={font=\scriptsize\sffamily}]
  % \panel{x0}{child page-1 frame}{copy?}
  \def\panel#1#2{%
    \begin{scope}[shift={(#1,0)}]
      \node[hd] at (6,37.5) {\iflangja{親}{parent}};
      \node[hd] at (28,37.5) {\iflangja{メモリ}{memory}};
      \node[hd] at (50,37.5) {\iflangja{子}{child}};
      % frames 0..4, height 7, centres 3.5 + 7 i
      \foreach \i in {0,...,4} { \draw[oth] (22,7*\i) rectangle ++(12,7); }
      \foreach \i/\n in {3/0, 1/1, 4/2} {
        \draw[ro] (22,7*\i) rectangle ++(12,7) node[midway, lab] {\iflangja{ページ}{page} \n};
      }
      % parent pages 0..2 (centres 10, 18, 26)
      \foreach \j/\i in {0/3, 1/1, 2/4} {
        \draw[ro] (0,6+8*\j) rectangle ++(12,8) node[midway, lab] {\j};
        \draw[->] (12,10+8*\j) -- (22,3.5+7*\i);
      }
      #2
    \end{scope}}
  % (a) after fork: child maps the same frames
  \panel{0}{%
    \foreach \j/\i in {0/3, 1/1, 2/4} {
      \draw[ro] (44,6+8*\j) rectangle ++(12,8) node[midway, lab] {\j};
      \draw[->] (44,10+8*\j) -- (34,3.5+7*\i);
    }}
  \node[capt] at (28,-5) {\iflangja{(a) \texttt{fork} 直後}{(a) after \texttt{fork}}};
  % (b) after the child writes page 1
  \panel{62}{%
    \foreach \j/\i in {0/3, 2/4} {
      \draw[ro] (44,6+8*\j) rectangle ++(12,8) node[midway, lab] {\j};
      \draw[->] (44,10+8*\j) -- (34,3.5+7*\i);
    }
    \draw[rw] (44,14) rectangle ++(12,8) node[midway, lab] {1};
    \draw[rw] (22,0) rectangle ++(12,7) node[midway, lab] {\iflangja{コピー}{copy}};
    \draw[->, very thick, ln-accent] (44,18) -- (34,3.5);}
  \node[capt] at (90,-5) {\iflangja{(b) 子がページ 1 に書き込んだ後}{(b) after the child writes page 1}};
  % legend
  \draw[ro] (0,-12) rectangle ++(4,3);
  \node[lab, anchor=west] at (5,-10.5) {\iflangja{書き込み禁止}{write-protected}};
  \draw[rw] (32,-12) rectangle ++(4,3);
  \node[lab, anchor=west] at (37,-10.5) {\iflangja{書き込み可能}{writable}};
\end{tikzpicture}

  \caption{コピーオンライト．(a) \texttt{fork} の後，親と子は同じ書き込み
    禁止のページフレームを対応付ける．(b) 子のページ 1 への書き込みが
    フォールトし，オペレーティングシステムはそのページだけをコピーする．}
  \label{fig:l08-cow}
\end{figure}

\begin{example}
  \SI{4}{\kibi\byte} のページ 25\,600 個（\SI{100}{\mebi\byte}）のアドレス
  空間を持つプロセスが \texttt{fork} し，子は \texttt{exec} を呼ぶ前に
  スタックとデータの 40 ページに書き込むとする．コピーオンライトは 25\,600
  ページではなく 40 ページ，すなわちアドレス空間の \SI{0.16}{\percent} を
  コピーする．
\end{example}

\section{チェックポインティング}
\label{sec:l08-checkpointing}

同じハードウェアによる支援は，障害管理にも役立つ．\source{P3L2，障害管理：
チェックポインティング．}

\begin{definition}[チェックポインティング]
  プロセスの状態を定期的に保存する障害・回復管理の技法を
  \term[ちぇっくぽいんてぃんぐ]{チェックポインティング}[checkpointing]と呼ぶ．
  障害の後，プロセスは最初からではなく，最新のチェックポイントから再開される．
\end{definition}

障害は避けられないかもしれないが，チェックポイントがあれば回復ははるかに
速い．最も単純な方法は，プロセスを停止して状態全体をコピーすることである．
大きなアドレス空間では停止が長くなり，繰り返すなら代価も高い．MMU はより
良い方法を可能にする．
\begin{itemize}
  \item オペレーティングシステムはプロセスのすべてのページを書き込み禁止に
    し，プロセスを走らせたまま状態全体を一度コピーする．
  \item それ以降は，書き込みフォールトまたはダーティビットによって，プロセス
    がどのページを変更したかを記録し，後のチェックポイントはこれらの汚れた
    ページだけをコピーする．このような\emph{増分}チェックポイントは，前の
    チェックポイントとの差分を格納する．
  \item プロセスを復元するには，完全なコピーとそれに続く差分から状態を
    再構成する．差分はバックグラウンドで併合しておくこともできる．
\end{itemize}

チェックポインティングは，さらなるサービスの基礎となる．
\begin{description}
  \item[デバッグ] \emph{巻き戻し再実行}では，プログラムをチェックポイント
    から再開（巻き戻し）して実行し直す．誤りが再現しなければ，見つかるまで
    順により古いチェックポイントを使う．
  \item[マイグレーション]
    \term[ぷろせすまいぐれーしょん]{プロセスマイグレーション}[process migration]では，
    プロセスのチェックポイントを別のマシンへ転送し，そこでプロセスを続行
    する．マイグレーションは，災害復旧や，プロセスをより少ないマシンへ集約
    するのに役立つ．\margin{Clark ら \cite{clark2005} は，
    実行中の Xen の仮想マシンを，ウェブサーバで \SI{210}{\milli\second}，
    ゲームサーバで \SI{60}{\milli\second} の停止時間で移送した．}
\end{description}
マイグレーションの停止時間を短く保つため，チェックポイントは速い繰り返しで
行われる．オペレーティングシステムは，プロセスを走らせたまま全ページを
コピーし，次にそのコピーの間に汚れたページを，次に2回目のコピーの間に汚れた
ページを，というように続ける．残るページが十分に少なくなって最後に停止して
コピーすることが許容できるようになるまで---あるいは，プロセスがコピーと同じ
速さでページを汚すために停止が避けられなくなるまで---これを繰り返
す．\margin{この繰り返しは\emph{プリコピー}マイグレーションと呼ばれ
る．\emph{ポストコピー}はその逆で，まず移送先でプロセスを再開し，
ページはネットワーク越しにフォールトで取り寄せる．}

\needspace{14\baselineskip}
\begin{example}
  50\,000 ページのプロセスを，毎秒 25\,000 ページをコピーする接続を通じて
  マイグレーションする．その間プロセスは毎秒 2000 個の異なるページを汚す．
  プロセスを停止してすべてをコピーすると \SI{2}{\second} かかる．各ラウンドが
  前のラウンドの間に汚れたページをコピーする繰り返しでは，どのラウンドも
  直前のラウンドの \SI{8}{\percent} のページをコピーし，3 ラウンドの後には
  最後の停止は約 \SI{1}{\milli\second} で済む．
  \begin{center}
    \small
    \begin{tabular}{@{}lrrr@{}}
      \toprule
      ラウンド & コピーしたページ数 & 所要時間 & その間に汚れたページ数 \\
      \midrule
      1（実行中） & 50\,000 & \SI{2}{\second}          & 4000 \\
      2（実行中） & 4000    & \SI{160}{\milli\second}  & 320 \\
      3（実行中） & 320     & \SI{12.8}{\milli\second} & 26 \\
      最後（停止） & 26     & \SI{1.04}{\milli\second} & --- \\
      \bottomrule
    \end{tabular}
  \end{center}
\end{example}

\needspace{14\baselineskip}

\begin{keypoint}[章のまとめ]
  仮想メモリはアドレス空間を物理メモリから切り離す．オペレーティングシステム
  はメモリを割り当て，すべてのアクセスを調停する．そのためにページングまたは
  セグメンテーションと，アドレスを変換しフォールトを起こす MMU とを用いる．
  ページテーブルは VPN を PFN に対応付け（\cref{eq:l08-translation}），その
  エントリは有効，ダーティ，アクセス，保護の各ビットを持つ．ページフォールト
  は，MMU が完了できないアクセスをオペレーティングシステムに委ねる．平坦な
  ページテーブルは大きすぎるので（\cref{eq:l08-pt-size}），多段にするか，
  逆ページテーブルとハッシュ法を用いる．そして TLB が変換を高速に保つ
  （\cref{eq:l08-eat}）．大きなページは，内部断片化を代価として，ページ
  テーブルを小さくし TLB の到達範囲を広げる．アロケータは外部断片化と戦い，
  Linux はページフレームにバディアロケータを，カーネルオブジェクトにスラブ
  アロケータを用いる．デマンドページングはページフォールトの際にページを
  スワップインし，ページアウトデーモンはセカンドチャンスアルゴリズムのような
  LRU の近似でページを選んでスワップアウトする．MMU はさらに，コピーオン
  ライトと増分チェックポインティングを可能にし，後者はデバッグやプロセス
  マイグレーションを支える．
\end{keypoint}

\end{document}
